\documentclass[12pt,a4paper]{ctexbook}

% Packages
\usepackage{amsmath,amssymb,amsthm}
\usepackage{mathrsfs}
\usepackage{geometry}
\geometry{left=2.5cm,right=2.5cm,top=2.5cm,bottom=2.5cm}
\usepackage{fancyhdr}
\pagestyle{fancy}
\fancyhf{}
\fancyhead[LE,RO]{\thepage}
\fancyhead[RE]{\leftmark}
\fancyhead[LO]{\rightmark}

% Theorem environments
\newtheorem{axiom}{公理}[section]
\newtheorem{theorem}{定理}[section]
\newtheorem{lemma}{引理}[section]
\newtheorem{corollary}{推论}[section]
\theoremstyle{definition}
\newtheorem{definition}{定义}[section]
\theoremstyle{remark}
\newtheorem{exercise}{习题}[section]
\newtheorem{example}{例}[section]

\title{朴素集合论}
\author{Paul R. Halmos}
\date{}

\begin{document}
\maketitle
\tableofcontents

\chapter{外延公理}

一群狼、一串葡萄或一群鸽子都是事物集合的例子。集合的数学概念可以作为所有已知数学的基础。这本小册子的目的是阐述集合的基本性质。顺便提一下，为了避免术语单调，我们有时会用“集”来代替“集合”。“类”这个词也用在这个语境中，但这样做有一点风险。原因是在集合论的某些处理方式中，“类”具有特殊的技术含义。稍后我们会再次提及这一点。

这一展开过程不包括对集合的定义。这种情况类似于人们熟悉的初等几何公理化方法。该方法不提供点和线的定义，而是描述人们能用这些对象做些什么。这里采用的半公理化观点假定读者对集合是什么具有普通的、人的、直觉的（且常常是错误的）理解；阐述的目的是勾画出人们可以用它们正确做出的许多事情中的一部分。

集合，正如通常所理解的那样，有元素或成员。集合的一个元素可以是一头狼、一颗葡萄或一只鸽子。重要的是要知道，一个集合本身也可以是另一个集合的元素。数学中充满了集合的集合的例子。例如，一条线是点的集合；

平面上所有直线的集合是集合的集合（点的集合）的一个自然例子。令人惊讶的可能不在于集合可以作为元素出现，而在于对于数学目的来说，永远不需要考虑其他任何元素。在本书中，我们尤其要研究集合，集合的集合，以及类似的高度和复杂性有时令人恐惧的高塔——别无他物。通过举例，我们偶尔会谈到$\dots $的集合

白菜、国王等等，但此类用法始终应被解释为仅仅是一个有启发性的比喻，而非正在发展的理论的一部分。

集合论的主要概念，在完全公理化的研究中是主要原始（未定义）概念，即属于概念。如果 x 属于 A（x 是 A 的元素，x 包含于 A），我们写作

xeA.

这种希腊字母 epsilon 的变体如此频繁地用于表示“属于”，以至于将其用于其他任何用途几乎都被禁止。大多数作者将 $\in$ 永远限定于集合论用途，而当他们需要希腊字母表的第五个字母时，则使用 $\varepsilon$。

也许简短地离题谈谈集合论中的字母使用规范会有所帮助。在前一段中，使用小写字母和大写字母并没有强制性的理由；我们本可以写下，并且经常也会写下像 xey 和 AeB 这样的内容。然而，只要可能，我们将通过一种约定来非正式地表明集合在特定层级中的地位：字母表开头的字母表示元素，字母表末尾的字母表示包含它们的集合；类似地，相对简单的字体表示元素，较大且花哨的字体表示包含它们的集合。例如：

x eA,AeX,Xee.

集合之间的一种可能关系，比属于关系更基本的，就是相等关系。两个集合A和B的相等关系普遍用熟知的符号表示。

A=B;

A 与 B 不相等这一事实用 A$\neq$B 表示。

属于的最基本性质是它与相等的关系，这可以表述如下。

外延公理。两个集合相等当且仅当它们拥有相同的元素。

元素。

以更大的矫饰和更少的清晰度而言：集合由其外延确定。

认识到外延公理并不仅仅是相等性在逻辑上所必需的性质，而是一个关于属于关系的非平凡断言，这很有价值。理解这一点的一种方式是考虑一个部分类似的情况，其中外延公理的类似物

并不成立。例如，假设我们考虑的是人而不是集合，并且，若x和A是人，当x是A的祖先时，我们记作xeA。（一个人的祖先是他的父母、他父母的父母、他们的父母，等等。）外延公理的类比在这里会说，如果两个人相等，那么他们有相同的祖先（这是“仅当”部分，且为真），并且如果两个人有相同的祖先，那么他们相等（这是“当”部分，且为假）。

如果 A 和 B 是集合，并且若 A 的每个元素都是 B 的元素，我们称 A 是 B 的子集，或 B 包含 A，记作

AC B

或

B A.

该定义的措辞意味着每个集合都必须被认为包含于自身（A C A）；这一事实被描述为集合包含关系是自反的。（注意，对于这个词的同样意义，相等关系也是自反的。）如果A和B是满足A C B且A$\neq$B的集合，则使用“真”这个词（真子集，真包含）。如果A、B和C是满足A C B且B C C的集合，那么A C C；这一事实被描述为集合包含关系是传递的。（这个性质也是相等关系所具有的。）

如果 A 和 B 是满足 A $\subset$ B 且 B $\subset$ A 的集合，那么 A 和 B 有相同的元素，因此，根据外延公理，A = B。这一事实被描述为集合的包含关系是反对称的。（在这方面，集合的包含关系与相等性不同。相等性是对称的，即如果 A = B，那么必然有 B = A。）事实上，外延公理可以用这些术语重新表述：如果 A 和 B 是集合，那么 A = B 的充分必要条件是 A $\subset$ B 且 B $\subset$ A。相应地，几乎所有的两个集合 A 和 B 相等的证明都分为两部分：首先证明 A $\subset$ B，然后证明 B $\subset$ A。

注意，属于（e）和包含（C）在概念上确实是完全不同的事物。一个重要的区别已经在上文中显现出来：包含总是自反的，而属于是否自反则完全不清楚。即：AC A 总是真的；AeA 曾经是真的吗？对于任何人所见过的任何合理集合，这肯定都不是真的。同样地，注意包含是传递的，而属于则不是。日常的例子，比如其成员本身也是组织的超级组织，对于感兴趣的读者来说会很容易想到。


\chapter{分离公理}

集合论的所有基本原理，除了外延公理之外，都是为了从旧集合构造新集合。这些构造集合的基本原理中的第一条且最重要的一条，粗略地说，就是：对于集合的元素，任何可以明确断言的性质都确定了一个子集，即该断言为真的那些元素构成的子集。

在精确表述这一原则之前，我们先看一个启发性的例子。设 A 是所有男人的集合。句子“x 是已婚的”对于 A 中的某些元素 x 成立，而对另一些则不成立。我们正在说明的原则，正是那个允许从给定的集合 A 过渡到由给定句子所指定的子集（即所有已婚男人的集合）的原则。为了指示该子集的生成，通常记作

$\{$ x $\in$ A : x 已婚 $\}$。

类似地

{xeA: x 未婚}

是全体单身汉的集合；

{xe A:x的父亲是亚当}

是包含该隐和亚伯且别无其他元素的集合；并且

$\{$x $\in$ A：x 是亚伯的父亲$\}$

是只包含 Adam 且不包含其他任何元素的集合。警告：一个装有帽子且只装有帽子的盒子与帽子不是一回事，同样，在此示例列表中的最后一个集合不应与

亚当。集合与盒子之间的类比有许多不足之处，但有时它能提供一个理解事实的有用图景。

支撑上述例子的精确一般表述所缺少的，仅仅是语句的定义。这里给出一个快速且非正式的定义。语句有两种基本类型，即，归属断言，

以及等式断言，

所有其他句子都通过重复应用通常的逻辑算子从这些原子句子得到，仅需遵守最低限度的语法和无歧义性。为了使定义更加明确（且更长），有必要在它后面附加一个“通常的逻辑算子”列表及句法规则。一个充分（且实际上冗余）的逻辑算子列表包含七个项目：

和，

或（即“要么$\dots $要么$\dots $或两者都”之意），非，

若—则—（或 蕴含），当且仅当，

有的（或存在），任意

关于句子构造的规则，可以描述如下。 ( i ) 将“not”置于一个句子之前，并将结果置于括号内。（这里以及下文中使用括号的理由是为了确保无歧义性。顺便提及，这使得所有其他标点符号都变得多余。句子定义所要求的完整括号体系很少全部需要。我们总是会尽量省略那些看起来不危及无歧义性的括号。在正常的数学实践中——本书也将遵循——会使用几种不同大小和形状的括号，但这仅是为了便于视觉辨识而已。） ( ii ) 在两个句子之间放置“and”或“or”或“if and only if”，并将结果置于括号内。 ( iii ) 将“if—then —”中的破折号替换为句子，并将结果置于括号内。 ( iv ) 将“for some—”或“for all—”中的破折号替换为一个字母，然后紧接一个句子，并将整个结构置于括号内。（如果所用字母未在句子中出现，也不会有任何问题。按照惯常且自然的约定，“for some y (x $\in$ A$)$”恰好意味着“x $\in$ A”。这同样）

如果所使用的字母已经用于“对于某些——”或“对于所有——”，则无害。回想一下，“对于某个 x (x$\in$A$)$”与“对于某个 y (y$\in$A$)$”意思相同；由此可见，明智地改变符号总能避免字母冲突。）

我们现在准备阐述集合论的主要原理，通常以其德文名称 Aussonderungsaxciom 来指称。

分类公理。对每一个集合A和每一个条件S(x)，都对应一个集合B，其元素恰好是A中那些使得S(x)成立的元素x。

“条件”在这里只是一个语句。符号表示意在表明字母x在语句S(x)中是自由的；这意味着x在S(x)中至少出现一次，且不是由“对某个x”或“对所有x”这类短语引入的。作为外延公理的直接推论，分类公理唯一地确定了集合B。为了表示B是由A和S(x)得到的方式，通常写作

B={xeA:S(x)}.

为了得到分类公理的一个有趣而富有启发性的应用，考虑作为S(x)的语句

非 (x ex).

在这里以及全文中，为了方便，我们将写作“xe A”（也写作“xA”）而不是“不属于”；在此记号下，S$($x$)$的角色现在由$\dots $扮演

xe'x.

由此可得，无论集合A是什么，如果B=$\{$x$\in$A: x$\notin$x$\}$，那么，对于所有的y，

(*) yeB 当且仅当 (yeA 且 ye'y)。

情形是否可能 \(BeA\)？我们着手证明答案是否定的。事实上，如果 \(BeA\)，那么要么也有 \(BeB\)（不太可能，但并非明显不可能），要么 \(Be'B\)。如果 \(BeB\)，那么根据(*)，假设 \(BeA\) 导致 \(Be'B\)——矛盾。如果 \(B \notin' B\)，那么再次根据(*)，假设 \(BeA\) 导致 \(BeB\)——又矛盾。这就完成了 \(BeA\) 不可能的证明，因此我们必有 \(BeA\)。这个结论最有趣的部分是，存在着某个不属于 \(A\) 的东西（即 \(B\)）。此论证中的集合 \(A\) 是相当任意的。换言之，我们证明了

无包含一切

或者，更壮观的是，

不存在全集。

这里“Universe”一词是在“论域”的意义上使用的，即在任何特定的讨论中，一个包含所有进入该讨论的对象的集合。

在较老的（前公理化）集合论方法中，全集的存在被视为理所当然，而前一段的论证被称为罗素悖论。其教训是，尤其在数学中，不可能无中生有。要指定一个集合，仅仅念出一些魔法词语（它们可能构成诸如“x$\in$'x”的句子）是不够的；还必须手头有一个集合，其元素适用于这些魔法词语。


\chapter{无序对}

尽管我们至今所说的一切可能让我们像是在真空中运作。为使讨论具有实质内容，现在让我们正式假定

存在一个集合。

由于稍后我们将提出一个更深刻、更有用的存在性假设，因此这个假设只起到临时作用。这个看似无害的假设的一个推论是，存在一个完全没有任何元素的集合。实际上，如果A是一个集合，将分离公理应用于A，使用语句“x$\neq$x”（或者，就此而言，任意其他恒假的语句）。结果是集合 $\{$x$\in$A: x$\neq$x$\}$，而该集合显然没有元素。外延公理意味着只能有一个没有元素的集合。该集合通常用符号

该集合被称为空集。

空集是每个集合的子集，换言之，ø $\subset$ A 对每个 A 成立。要确立这一点，我们可以这样论证。要证明的是，ø 中的每个元素都属于 A；由于 ø 中没有元素，该条件自动满足。这个推理是正确的，但可能不那么令人满意。由于这是常见现象的一个典型例子——一个在“空洞”意义上成立的条件——给缺乏经验的读者一句忠告或许合适。要证明关于空集的某事为真，就证明它不可能为假。例如，ø $\subset$ A 怎么可能为假呢？只有当 ø 有一个元素不属于 A 时，它才可能为假。由于 ø 根本没有元素，这是荒谬的。结论：ø $\subset$ A 不为假，因此对每个 A 有 ø $\subset$ A。

迄今为止发展的集合论仍然相当贫乏；据我们所知，只有一个集合，且该集合是空集。是否有足够的集合来确保每个集合都是某个集合的元素？对任意两个集合，是否总存在第三个集合，使得它们都属于它？三个集合呢？四个呢？任意多个呢？我们需要一个新的集合构造原理来解决这类问题。下面的原理是一个好的开端。

配对公理。对任意两个集合，存在一个集合，它们都属于该集合。

请注意，这仅仅是对上述第二个问题的肯定回答。

为了安抚担忧者，让我们赶紧指出，诸如“two”这样的词

上面使用的“三”和“四”这两个词，并不是指那些后面将要定义的、同名的数学概念；目前这些词仅仅是“某物然后再某物”重复适当次数的普通语言缩写。因此，举例来说，配对公理在未缩写的形式下说的是：如果a和b是集合，那么存在一个集合A使得aeA且beA。

配对公理的一个推论（实际上是一个等价表述）是，对于任意两个集合，存在一个集合恰好包含它们且不包含其他任何元素。确实，如果 a 和 b 是集合，且如果 A 是一个集合使得 a $\in$ A 且 b $\in$ A，那么我们可以将分离公理应用于 A 和语句 “x = a 或 x = b。” 结果就是集合

{xeA: x=a 或 x =b}

而那个集合显然只包含a和b。外延公理意味着具有这个性质的集合只能有一个。那个集合的常用符号是

{a, b};

该集合被称为由 a 和 b 构成的对（或者，为了与后续概念进行强调对比，称为无序对）。

如果暂且将语句“x = a 或 x = b”记作 S(x)，那么我们可以这样表述配对公理：存在一个集合 B 使得

(*) xe B 当且仅当S(x).

分类公理，应用于一个集合 A，断言存在一个集合 B 使得

(**) zeB 当且仅当 (xeA 且 S(x))。

(*)与(**)之间的关系体现了一种相当常见的情形。所有其余的集合构造原理在(*)作为(**)的伪特例的意义下，都是分离公理的伪特例。它们都断言由某一条件所指定的集合的存在性；如果事先知道存在一个包含所有指定元素的集合，那么一个仅包含这些元素的集合的存在性确实可作为分离公理的一个特例而推出。

如果 a 是一个集合，我们可以构造无序对 {a,a}。该无序对记作

并称 {a} 为 a 的单元素集；它被唯一地描述为仅以 a 为其元素。因此，例如， 和 {Ø} 是非常不同的集合；前者没有元素，而后者有唯一元素 ø。说 aeA 等价于说 {a}CA。

配对公理确保每个集合都是某个集合的元素，并且任意两个集合同时是某个同一集合的元素。（关于三个、四个及更多集合的相应问题将在后面回答。）另一个相关的评论是，从我们迄今所做的假设中，我们确实可以推断出非常多集合的存在。例如，考虑集合 ø,{x},{{Ø}},{{Ø}}} 等；考虑由它们中任意两个构成的配对，如 {ø,{x}}；考虑由任意两个这样的配对构成的配对，或者由任意单元素集和任意配对构成的混合配对；并如此无限进行下去。

习题。用这种方法得到的所有集合是否互不相同？

在继续我们关于集合论的学习之前，我们暂停片刻讨论一个记号问题。用记号表示所描述的集合 B 似乎很自然

在(*)中用{x:S(x)}表示；在曾考虑过的特殊情形中，{x:x = a 或 x = b} = {a,b}。

我们将在方便且允许的情况下使用这种符号表示法。也就是说，如果S(x)是一个关于x的条件，使得S(x)所指定的那些x...

fies 构成一个集合，那么我们可以将该集合记为 {x:S(x)}。

在 A 是集合且 S$($x$)$ 为 (x$\in$A$)$ 的情况下，那么构造 $\{$x:S$($x$)\}$ 是允许的；事实上

{x:xe A}=A.

如果 A 是一个集合且 S(x)是任意一个语句，那么可以构成 {x:xeA and S(x)}；这个集合与 {xeA:S(x)} 相同。作为进一步的例子，我们注意到

集合 $\{$z : x $\neq$ x$\}$ 等于空集

和

当 S$($x$)$ 为 (x $\in$ x$)$ 时，或当 S$($x$)$ 为 (x = x$)$ 时，所述的那些 x 并不构成一个集合。

尽管有不能无中生有的格言，被告知某些集合并非真正的集合，甚至它们的名字都绝不能提及，这似乎有点严厉。有些集合论方法试图通过系统地使用这种非法集合，只是不称其为集合，来减轻打击；习惯用的词是“类”。精确解释类究竟是什么以及如何使用它们，在当前的方法中并不重要。粗略地说，一个类可以等同于一个条件（语句），或者更确切地说，等同于一个条件的“外延”。


\chapter{并集与交集}

若A与B是集合，有时自然地希望将其元素合并成一个综合集合。描述这种综合集合的一种方式是要求它包含至少属于对{A,B}中两个成员之一的所有元素。这一表述暗示了对其自身的广泛推广；类似的构造当然应适用于任意集族，而不仅仅是集合对。换言之，所需要的是如下集合构造原理。

并集公理：对于任意集合族，存在一个集合，它包含了所有至少属于该族中某一个集合的元素。

再次说明如下：对于每一个集族 e，存在一个集合 U，使得如果对于 e 中的某个 X 有 x e X，那么 x e U。（注意“至少一个”与“某个”相同。）

上文描述的广泛集合U可能过于宽泛；它可能包含不属于集合族e中任何集合X的元素。这很容易补救；只需应用分类公理来形成

集合{xe U:xeX 对于某个 X 属于 e}。

集合

（此处的条件是转换为惯用语，即数学上更易接受的“for some X(xeXand Xee)”。）由此可得，对于每个 x，x 属于该集合的充要条件是，存在 e 中的某个 X 使得 x 属于 X。如果改换记法并再次称新集合为 U，那么

U=$\{$x: x$\in$X 对某个 X$\in$e$\}$

这个集合 U 被称为集合族 e 的并集；注意

外延公理保证其唯一性。U的最简符号

根本被使用的那个在数学界并不很流行；它是

大多数数学家更喜欢像这样的

U {X:Xee}

函数 \( e^{-x^{2}} \) 没有可以用初等函数表示的原函数，但我们仍然可以计算它在整个实数轴上的定积分。技巧是通过转换到极坐标来计算该积分的平方： \[ \left(\int\_{-\infty}^{\infty} e^{-x^{2}} dx\right)^{2} = \int\_{-\infty}^{\infty} e^{-x^{2}} dx \int\_{-\infty}^{\infty} e^{-y^{2}} dy = \int\_{-\infty}^{\infty} \int\_{-\infty}^{\infty} e^{-(x^{2}+y^{2})} dx dy = \int\_{0}^{2\pi} \int\_{0}^{\infty} e^{-r^{2}} r dr d\theta = \pi. \] 因此， \[ \int\_{-\infty}^{\infty} e^{-x^{2}} dx = \sqrt{\pi}. \] 这是概率论与统计学的一个基本结果。

Ux.eX.

在某些重要的特殊情况下，还有其他替代方案可用；它们将在适当的时候予以说明。

目前，我们将对并集理论的研究仅限于最简单的事实。所有事实中最简单的是

U{X:XeØ}=Ø

下一个最简单的事实是

U {X:Xe{A}}=A.

在上述极为简洁的记号中，这些事实表示为

Uø=x

和

U{A}=A

证明由定义直接可得。

在集合对的并集中有更多的实质内容（这毕竟就是引发整个讨论的原因）。在这种情况下使用了特殊符号：

U{X:Xe{A,B}}=AU B

并集的一般定义在这种特殊情况下意味着 xeA U B

当且仅当x属于A或B或两者；由此可得A $\cup$ B = $\{$x : x $\in$ A 或 x $\in$ B$\}$。

以下是一些关于对集的并集的容易证明的事实：

A Uø = A

AU B=BU A（交换律）,

A U(BU C)=(A U B)UC(结合律)，

AU A=A（幂等性），

ACB 当且仅当 A U B = B。

每个数学学生一生中至少应该亲自证明一次这些事情。这些证明基于逻辑运算符“或”的相应基本性质。

一个同样简单但颇具启发性的事实是 {a}U {b}={a,b}。

这提示了推广有序对的方法。具体而言，我们写

{a,b,c}={a}U{b}U {c}.

该方程定义了其左侧。其右侧按理说应至少有一对括号，但鉴于结合律，省略它们不会引起误解。由于容易证明

{a,b,c}={x:x =a 或 x =b 或 x =c}

现在我们知道，对于任意三个集合，存在一个集合包含它们且仅包含它们；自然称那个唯一确定的集合为由它们构成的（无序）三元组。如此引入的记号和术语到更多项（四元组等）的推广是显然的。

并集的构造与另一种集合论运算有许多相似之处。如果 A 和 B 是集合，A 与 B 的交集是集合

由$\dots $定义

A$\cap$B=$\{$xeA:xeB$\}$.

该定义在A和B中是对称的，即使看起来并非如此；我们有A$\cap$B=$\{$xeB:xeA$\}$,

并且，事实上，由于 x $\in$ A$\cap$B 当且仅当 x 既属于 A 又属于 B，由此可知

A$\cap$B=$\{$x:xe A 且 xeB$\}$.

关于交集的基本事实，以及它们的证明，类似于关于并集的基本事实：

A$\cap$ø=x,

A$\cap$B=B$\cap$A,

A$\cap($B$\cap$C$)$=$($A$\cap$B$)\cap$C

A$\cap$A=A

ACB 当且仅当 A$\cap$B=A。

交集为空集的集合对经常出现，因此有必要引入一个专门的术语：如果 A$\cap$B=$\varnothing$，则称集合 A 和 B 不交。该术语有时也用于一个集合族，表示该族中任意两个不同的集合都不交；或者，在这种情况下我们也可以称该族为两两不交族。

关于并集和交集的两个有用事实同时涉及这两种运算：

A$\cap($B UC$)$=$($A$\cap$B$)$U$($A$\cap$C$)$,

A U$($B$\cap$C$)$=$($A U B$)\cap($A U C$)$.

这些恒等式称为分配律。作为集合论证明的一个示例，我们证明第二个。如果x属于左边，那么x要么属于A，要么既属于B又属于C；如果x在A中，那么x既属于A U B又属于A U C，而如果x既在B中又在C中，那么同样，x也既属于A U B又属于A U C；由此可得，在任何情况下，x都属于右边。这证明了右边包含左边。为了证明相反的包含关系，只需注意到，如果x既属于A U B又属于A U C，那么x要么属于A，要么既属于B又属于C。

两个集合 A 和 B 的交的形成，或者不妨说，对集合对 $\{$A, B$\}$ 取交，是一种更一般的运算的特例。（这是交集理论模仿并集理论的另一个方面。）一般交运算的存在依赖于这样一个事实：对于每个非空的集合族，都存在一个集合恰好包含那些属于该给定族中每个集合的元素。换句话说：对于每个异于 ø 的集族 e，存在一个集合 V，使得 x$\in$V 当且仅当对于 e 中的每个 X 都有 x$\in$X。要证明这个断言，令 A 为 e 中任意特定的集合（这一步由 e$\neq$x 的事实保证），并写出

V={xeA:xeX 对于 e 中的每个 X}

（该条件意味着“对所有X（如果Xee，那么xeX）。”）这种依赖

V对A的任意选择的依赖性是虚幻的；事实上

V={x:对于 e 中的每一个 X，xeX}。

集合 V 被称为集合族 e 的交集；外延公理保证了它的唯一性。通常的记号类似于并集的记号：不是采用那种无可反对但却不受欢迎的

集合 V 通常表示为

n{X:Xee}

或者

$\cap$x.eX.

EXERCISE. (A$\cap$B$)$ U C = A$\cap($B U C$)$ 的一个充要条件是 CCA。注意到该条件与集合 B 无关。


\chapter{补集与幂集}

如果A和B是集合，A与B的差，更常被称为

作为 B 在 A 中的相对补集，集合 A-B 定义为 A-B={xeA: xeB}。

注意，在这个定义中，不需要假定 BC A。为了尽可能简单地记录关于取补运算的基本事实，我们仍然假定（仅限于本节）所有要提及的集合都是同一个集合 E 的子集，并且所有补集（除非另有说明）都是相对于该 E 来形成的。在这种情况下（这种情况相当常见），记住基础集合 E 比一直写下它更容易，这就使得简化符号成为可能。一个经常用来表示 A 的暂时绝对（相对于相对而言）补集的符号是 A'。用这个符号，关于取补运算的基本事实可以陈述如下：

(A'$)$'=A, $\varnothing$′=E, E′=$\varnothing$

A$\cap$A'=x，AU A′=E，

A $\subseteq$ B 当且仅当 B' $\subseteq$ A'。

关于补集的最重要的定律是所谓的德摩根定律：

(A U B$)$'=A'$\cap$B',$($A$\cap$B$)$'=A'U B'.

（我们很快就会看到，德摩根律对于比仅仅成对集合更大的集族的并集和交集也成立。）这些关于

补运算意味着集合论中的定理通常成对出现。如果在涉及E的子集的并、交和补的包含关系或等式中，我们将每个集合替换为其补集，交换并与交，并反转所有包含关系，则结果将是另一个定理。这一事实有时被称为集合的对偶原理。

这里有一些关于补集的简单练习题。

A-B=A$\cap$B'.

ACB当且仅当A-B=ø。

A-$($A-B$)$=A$\cap$B

A$\cap($B-C$)$=$($A$\cap$B$)$-$($A$\cap$C$)$.

A$\cap$BC$($A$\cap$C$)$U$($B$\cap$C'$)$

(A U C$)\cap($BU C"$)$CAU B.

如果 A 和 B 是集合，A 与 B 的对称差（或布尔和）是集合 A+B，定义为

A+B=(A-B)U(B-A)

此运算满足交换律$($A+B=B+A$)$和结合律$($A+ (B+C$)$=$($A+B$)$+C$)$，且满足A+$\varnothing$=A以及A+A

这可能正是澄清交集理论中一个琐细但偶尔令人困惑的部分的好时机。首先，回顾一下，交集仅对非空集合族有定义。理由是，相同的方法对空集合族未定义一个集合。句子中指定了哪些 x？

对于Ø中的每一个X，xe X？

关于 ø 的问题，通常通过相应的否定问题更容易看出答案。哪些 x 不满足所述条件？如果并非对 ø 中的每个 X 都有 xeX，那么必定存在 ø 中的一个 X 使得 xeX；然而，由于 Ø 中根本没有任何 X，这是荒谬的。结论：没有 x 不满足所述条件，或者等价地说，每个 x 都满足它。换句话说，条件所指定的 x 穷尽了（不存在的）全域。这里没有深奥的问题；仅仅是令人困扰地总是不得不做出

由于构造过程中某处的某个集合可能结果为空集，而不得不做出的限定和例外处理。对此我们无能为力；这不过是生活中的一个事实。

如果我们把注意力限制在某个特定集合E的子集上，正如我们临时约定那样，那么在前文所述的不愉快情况

ceding段落似乎消失了。关键是在那种情况下我们

可以定义集合族e（E的子集）的交集为该集合

{xeE:xe X 对于 e 中的每个 X}

这并非革命性的；对于每个非空集合族，新定义与旧定义一致。区别在于旧定义和新定义处理空集合族的方式；根据新定义，$\cap$x.øX 等于 E。（对于 E 中的哪些元素 x，会使得“对于 ø 中的所有 X，x$\in$X”不成立？）这种区别只是语言上的问题。稍加思考便会发现，对 E 的子集族 e 的交集所给出的“新”定义，实际上与旧定义下族 e$\cup\{$E$\}$ 的交集相同，而后者永不为空。

我们一直在考虑集合 E 的子集；这些子集自身能否构成一个集合？下面的原理保证了答案是肯定的。

幂集公理：对于每个集合，都存在一个集合，它的元素中包含该集合的所有子集。

换句话说，如果 E 是一个集合，那么存在一个集合（集族）P，使得若 X $\subseteq$ E，则 X $\in$ P。

前面描述的集合 o 可能比所需的大；它可能包含不是 E 的子集的元素。这很容易补救；只需应用分类公理来构造集合 $\{$X $\in$ o : X $\subset$ E$\}$。（回忆一下，“X $\subset$ E” 与 “对所有 x，若 x $\in$ X 则 x $\in$ E” 是一个意思。）由于对于每个 X，X 属于这个集合的充分必要条件是 X 是 E 的子集，因此，如果我们改换符号，再次称这个集合为 o，那么

①={X:XCE}.

集合 c 称为 E 的幂集；外延公理保证了其唯一性。o 对 E 的依赖性通过写 (E) 而不仅仅是 o 来表示。

因为集合 o(E) 与 E 相比非常大，所以不容易给出例子。如果 E=Ø，情况就足够清楚了；集合 o(x) 是

单元素集 {x}。单元素集和双元素集的幂集也很容易描述；我们有

0({a})={ø,{a}}

关键规则： 1. 数学符号（公式、符号、变量）必须保持不变 2. 集合论术语应遵循标准的中文数学惯例： - set = 集合 - element = 元素 - subset = 子集 - union = 并集 - intersection = 交集 - empty set = 空集 - power set = 幂集 - ordered pair = 有序对 - relation = 关系 - function = 函数 - domain = 定义域 - range = 值域 - cardinal number = 基数 - ordinal number = 序数 - axiom of choice = 选择公理 - well-ordering = 良序 - transfinite = 超限 3. 仅输出中文翻译，无解释，无注释 4. 保留段落结构和换行

①({a,b})={ø,{a},{b},{a,b}}.

三元集合的幂集有八个元素。读者或许可以猜出（并在此受到挑战去证明）包含所有这些陈述的推广：一个有限集合，比如有 n 个元素，其幂集有 2^n 个元素。（当然，像“有限”和“2^n”这样的概念对我们来说还没有正式的地位；但这不应妨碍我们对它们有非正式的理解。）n 作为指数（2 的 n 次方）出现，这与幂集之所以得名有关。

如果 e 是集合 E 的子集族（即，e 是...的子族）

(E)),则 写

①={Xe(E):X'ee}.

（为确保在定义 D 时所用的条件在精确的技术意义下是一个语句，它必须被改写成类似如下形式）

存在某个 Y[Yee 并且对于所有的 x(xeX 当且仅当 (xeE 且 xeY))]。

类似地，当我们希望使用定义的缩写而不仅仅是逻辑和集合论原语时，类似的评论也适用。这种翻译几乎不需要任何巧思，我们通常会省略它。）通常的做法是

用符号表示集族 D 的并集和交集

Ux.eX′和 $\cap$x.eX'.

在这种记法中，德摩根定律的一般形式变为

(Ux.eX$)$'=$\cap$x.eX′

以及 (i) 继续这一过程超越 \$\omega\$，构造 \$\mathcal{P}(E\_1), \mathcal{P}(E\_2)\$ 等，即通过沿着所有序数迭代应用幂集和并集运算，从而得到 *超限集合层级*； (ii) 陈述 *基础公理*，也称为 *正则公理*，它断言任何非空集合都包含一个与它不相交的元素，从而排除了集合自身作为自身元素的集合的存在，或更一般地，排除了无穷递减的 \$\in\$-链；以及 (iii) 最后，添加由弗兰克尔和斯科伦引入的 *替换公理*，它断言对于每个集合 \$x\$ 和每个由一阶公式定义的函数 \$f\$，像 \$f[x] = \{f(y) : y \in x\}\$ 也是一个集合。

($\cap$x.eX$)$'=Ux.eX'.

这些方程的证明是相应定义的直接结果。

练习。证明 o$($E$)\cap$o$($F$)$=$($E$\cap$F )且 o$($E$)$Uo$($F$)$C

(E U F$)$。这些断言可以推广到$\cap$x.eo$($X$)$=0$(\cap$x.X$)$

和

Ux.e(X)co(Ux.eX);

找到这些推广在此处表达时所用符号的一个合理解释，然后证明它们。进一步初等。

事实：

$\cap$x: (B$)$X=x,

和

如果 E CF，则 o(E)C0(F)。

一个有趣的问题涉及算子 o 和 U 的交换性。证明 E 总是等于 U o (E)（即 E = U o (E)），但以另一种顺序将 o 和 U 应用于 E 的结果是一个包含 E 作为子集的集合，通常是真子集。


\chapter{有序对}

将集合 A 的元素按某种顺序排列是什么意思？例如，假设集合 A 是由不同元素组成的四元组 {a,b,c,d}，并且假设我们想要按顺序考虑它的元素

cbd a。

即使没有关于这意味着什么的精确定义，我们也可以用它做一些集合论上的巧妙之事。我们可以，即，考虑，对于排序中的每个特定位置，所有出现在该位置或之前的元素的集合；我们以这种方式得到这些集合

{c} {c,b} {c,b,d} {c,b,d,a}

那么我们可以继续考虑集合（如果你觉得更顺耳，也可以叫收集）e={{a,b,c,d},{b,c},{b,c,d},{c}}

它恰好以那些集合作为它的元素。为了强调基于直观且可能模糊的顺序概念成功地产生了一些坚实而简单的东西，即一个朴素无修饰的集合e，上面以打乱的方式呈现了e的元素及其元素的元素。（倾向于字典序的读者或许能看出这种打乱方式中的某种方法。）

我们暂且继续假装我们确实知道序的含义。假设在匆匆一瞥前文时，我们只能抓住集合 e；我们能用它重新找回使之产生的序吗？答案显而易见是肯定的。检查 e 的元素（当然，它们本身都是集合），找出一个被所有其他元素包含的元素；由于 {c} 符合条件（且没有其他元素符合），我们知道 c 必定是第一个元素。接着寻找 e 的下一个最小元素，

即，在移除{c}后剩下的所有集合中都包含的那个；由于{b,c}符合要求（且没有其他集合符合），我们知道b必定是第二个元素。如此进行（只需再完成两步），我们便从集合e过渡到给定集合A的指定排序。

其寓意是：我们可能并不确切知道对集合A的元素进行排序意味着什么，但对于每一种排序，我们都可以关联一个由A的子集构成的集合e，使得原本的排序能从中唯一地复现出来。

e.（这里有一个非平凡的练习：找出那些对应A中某种次序的A的子集的集合的内在特征。由于“次序”对我们还没有正式的定义，整个问题在正式意义上是无意义的。后面的内容不依赖于这个问题的解答，但读者通过尝试寻找它会学到有价值的东西。）从A中一个次序到集合e再回来的过程，上文对于一个四元组已经举例说明；对于一对元素，一切至少简单两倍。若A={a,b}，且在所希望的次序中a最先出现，则e={{a},{a,b}}；然而若b最先出现，则e={{b},{a,b}}。

a 与 b 的有序对，其第一坐标为 a，第二坐标为 b，是集合 (a, b)，定义为

(a,b)={{a},{a,b}}.

无论这个定义的动机多么令人信服，我们仍须证明其结果具有有序对必须满足的主要性质，才配得上它的名称。我们必须证明：如果 (a,b$)$ 和 (x,y$)$ 是有序对，且 (a,b$)$ = (x,y$)$，那么 a=x 且 b=y。为了证明这一点，我们首先注意到，如果 a 和 b 恰好相等，那么有序对 (a,b$)$ 就是单元素集 $\{\{$a$\}\}$。反之，如果 (a,b$)$ 是单元素集，那么 $\{$a$\}$ = $\{$a,b$\}$，从而 b $\in\{$a$\}$，因此 a=b。现在假设 (a,b$)$ = (x,y$)$。如果 a=b，那么 (a,b$)$ 和 (x,y$)$ 都是单元素集，因此 x=y；由于 $\{$x$\}\in$ (a,b$)$ 且 $\{$a$\}\in$ (x,y$)$，可知 a、b、x、y 全都相等。如果 a$\neq$b，那么 (a,b$)$ 和 (x,y$)$ 都恰好包含一个单元素集，分别是 $\{$a$\}$ 和 $\{$x$\}$，所以 a=x。由于在这种情况下同样有 (a,b$)$ 和 (x,y$)$ 都恰好包含一个不是单元素集的无序对，分别是 $\{$a,b$\}$ 和 $\{$x,y$\}$，因此 $\{$a,b$\}$ = $\{$x,y$\}$，从而特别地有 b $\in\{$x,y$\}$。由于 b 不能等于 x（否则就会有 a=x 且 b=x，从而 a=b），我们必然有 b=y，证明完成。

如果 A 和 B 是集合，是否存在一个集合包含所有有序对 (a,b$)$，其中 a $\in$ A 且 b $\in$ B？很容易看出答案是肯定的。事实上，如果 a $\in$ A 且 b $\in$ B，那么 $\{$a$\}\subset$ A 且 $\{$b$\}\subset$ B，因此 $\{$a,b$\}\subset$ A $\cup$ B。又因为 $\{$a$\}\subset$ A $\cup$ B，所以两个 $\{$a$\}$

以及 $\{$a,b$\}$是 o$($AU B$)$的元素。这意味着$\{\{$a$\}$,$\{$a,b$\}\}$是 o$($A U B$)$的子集，因此它是 o$($0$($A U B$)$)的元素；换句话说，只要 aeA 和 beB，就有 (a,b$)$e$($C$($A U B$)$)。一旦已知这一点，就可以例行地应用分类公理和外延公理来产生唯一的集合 A$\times$B，它恰好由满足 a 在 A 中且 b 在 B 中的有序对 (a,b$)$ 组成。这个集合被称为

A 与 B 的笛卡尔积；它的特征是

A$\times$B=$\{$x:x=$($a,b$)$对于某个a在A中且对于某个b在B中$\}$。

两个集合的笛卡尔积是一个有序对集合（即，每个元素都是有序对的集合），并且笛卡尔积的每个子集也是如此。重要的是在技术层面上我们知道反过来也成立：每个有序对集合都是两个集合的笛卡尔积的子集。换句话说：如果 R 是一个集合，且 R 的每个元素都是有序对，那么存在两个集合 A 和 B 使得 RCA$\times$B。证明是初等的。假设确实 xeR，那么对于某个 a 和某个 b，x =$\{\{$a$\}$,$\{$a,b$\}\}$。问题在于从花括号下挖出 a 和 b。由于 R 的元素都是集合，我们可以构造 R 中集合的并集；因为 x 是 R 中的一个集合，所以 x 的元素属于该并集。由于 $\{$a,b$\}$ 是 x 的元素之一，我们可以用前面所谓的粗暴记法写出，$\{$a,b$\}$eUR。一对花括号消失了；让我们再做一次同样的事情，让另一对花括号也消失。构造 UR 中集合的并集。因为 $\{$a,b$\}$ 是这些集合之一，所以 $\{$a,b$\}$ 的元素属于那个并集，因此 a 和 b 都属于 UUR。这就实现了前面所说的；为了将 R 表示为某个 A$\times$B 的子集，我们可以取 A 和 B 都为 UUR。通常希望取 A 和 B 尽可能小。为此，只需应用分离公理来产生这些集合

A=$\{$a: 存在b使得$($a,b$)\in$R$\}$

和

B={b: 存在某个a使得((a,b)eR)}.

这些集合分别称为 R 在第一坐标和第二坐标上的投影。

无论集合论如今多么重要，在其创立之初，一些学者视其为一种疾病，并希望数学能很快从中康复。因此，许多集合论的考量被称为“病态的”，而这个词至今仍在数学中使用；它常常指代说话者不喜欢的东西。一个明确界的

有序对((a,b)={{a},{a,b}})常常被归入病态集合论。为了那些认为在这种情况下这个名称实至名归的人，我们指出，该定义现在已经完成了它的使命，今后将不再使用。我们需要知道的是，有序对由其第一坐标和第二坐标决定并唯一决定它们，笛卡尔积可以被构造，并且每个有序对的集合都是某个笛卡尔积的子集；至于用哪种具体的方法来实现这些目标则是无关紧要的。

要找出许多数学家对上述有序对的显式定义所持不信任与怀疑的根源并不难。问题并不在于有任何错误或遗漏；我们所定义的概念其相关性质全都是正确的（即符合直觉要求），所有正确性质也都具备。问题在于这个概念有一些无关的、偶然且分散注意力的性质。定理 (a, b$)$ = (x, y$)$ 当且仅当 a = x 且 b = y 是我们期望学到的关于有序对的知识。另一方面，$\{$a, b$\}\in$ (a, b$)$ 这一事实似乎是偶然的；这是定义的一个怪异性质，而非概念的内在性质。

人为性的指责是真实的，但为了概念的经济性，这并不是太高的代价。有序对的概念本可以作为附加的原始概念引入，并以公理方式赋予恰好恰当的性质，不多也不少。在一些理论中就是这样做的。数学家的选择在于要么记住几个额外的公理，要么遗忘几个偶然的事实；这种选择显然是一个品味问题。类似的选择在数学中频繁出现；例如，在本书中，我们将在各种数的定义中再次遇到它们。

练习. 如果 A、B、X 和 Y 是集合，那么

(i$)$(A U B$)\times$X=$($A$\times$X$)$U$($B$\times$X$)$

(ii$)$(A$\cap$B$)\times($XnY$)$=$($A$\times$X$)\cap($B$\times$Y$)$

iii$)$(A-B$)\times$X=$($A$\times$X$)$-$($B$\times$X$)$.

如果 A=x 或 B=x，那么 A$\times$B=x，反之亦然。如果

ACX 和 BCY，则 A$\times$BCX$\times$Y，并且（假设 A$\times$B$\neq$ ø）反之亦然。


\chapter{关系}

利用有序对，我们能够以集合论的语言来阐述关系的数学理论。这里所说的关系，是指诸如婚姻（男人与女人之间）或属于（元素与集合之间）之类的事物。更明确地说，我们称之为关系的，有时被称为二元关系。三元关系的一个例子是人与人之间的亲子关系。

（亚当和夏娃是该隐的父母）。在本书中，我们将没有机会讨论三元、四元或更复杂的关系理论。

观察任何具体的关系，例如婚姻，我们可能会考虑某些有序对 (x,y)，即那些 x 是男人、y 是女人且 x 与 y 结婚的有序对。我们尚未给出关系的一般概念的定义，但似乎合理的是，正如在这个婚姻例子中，每个关系都应唯一地确定所有那些第一个坐标与第二个坐标具有该关系的有序对所组成的集合。如果我们知道这个关系，我们就知道这个集合；并且更好的是，如果我们知道这个集合，我们就知道这个关系。例如，如果我们得到了与婚姻相对应的人的有序对集合，那么即使我们忘记了婚姻的定义，我们也总能判断男人 x 是否与女人 y 结婚；我们只需看有序对 (x,y) 是否属于该集合。

我们或许不知道关系是什么，但我们确实知道集合是什么，而前文的种种考虑建立了关系与集合之间的紧密联系。对关系的精确集合论处理正是利用了这种启发性的联系；最简单的做法就是将关系定义为相应的集合。这正是我们所做的；我们特此将关系定义为一个有序对的集合。明确地说：若集合 R 中的每个元素都是有序对，则它是一个关系。

R 的元素是一个有序对；这当然意味着，如果 z$\in$R，那么存在 x 和 y 使得 z=$($x,y$)$。如果 R 是一个关系，有时为方便起见，可以通过写法表示 (x,y$)\in$R 这一事实

zR y

并且，像日常语言那样，说 z 与 y 具有关系 R。

最无趣的关系是空关系。（要证明 ø 是有序对的集合，只需在 ø 中寻找一个不是有序对的元素。）另一个无趣的例子是任意两个集合 X 和 Y 的笛卡儿积。这里有一个稍有趣些的例子：设 X 为任意集合，并设 R 为 X$\times$X 中所有使得 x=y 的有序对 (x,y$)$ 的集合。关系 R 正是 X 中元素间的相等关系；若 x 和 y 属于 X，则 xRy 与 x=y 含义相同。眼下再举一个例子便足够了：设 X 为任意集合，并设 R 为 X$\times$ o$($X$)$ 中所有使得 xeA 的有序对 (x,A$)$ 的集合。这个关系 R 正是 X 中元素与 X 的子集之间的属于关系；若 xeX 且 Ae$($X$)$，则 x R A 与 x eA 含义相同。

在前一节中，我们看到了与每一个有序对集合 R 相关联的有两个集合，称为 R 在第一坐标和第二坐标上的投影。在关系理论中，这些集合被称为 R 的定义域和值域（缩写为 dom R 和 ran R）；我们回顾一下，它们定义为

dom R = {x: 存在 y (xRy)}

抱歉，您似乎未提供需要翻译的英文文本。请提供具体内容，我将严格按照数学专业翻译规范进行中文翻译。

ran R ={ y: 存在某个 x (x R y) }.

若 R 是婚姻关系，使得 ax Ry 意味着 x 是男人，y 是女人，且 x 和 y 彼此结婚，则 dom R 是已婚男人的集合，而 ran R 是已婚女人的集合。空集 ø 的定义域和值域都等于 ø。若 R=X$\times$Y，则 dom R=X

且 ran R = Y。若 R 是 X 中的相等关系，则 dom R = ran R = X。若 R 是属于关系，介于 X 与 o(X) 之间，则 dom R = X 且 ran R = o(X) \ {Ø}。

如果 R 是包含在笛卡儿积 X$\times$Y 中的一个关系（因此 dom R $\subset$ X 且 ran R $\subset$ Y），有时为了方便，我们说 R 是从 X 到 Y 的关系；若非从 X 到 X 的关系，我们可以说关系在 X 中。X 中的一个关系 R 是自反的，如果对于 X 中的每个 x 有 xRx；是对称的，如果 xRy 蕴含 yRx；是传递的，如果 xRy 和 y R z 蕴含 z R z。（练习：对于这三种可能的性质中的每一种，找出一个不具有该性质但具有其他两种性质的关系。）一个关系

集合中的一个关系是等价关系，如果它具有自反性、对称性和传递性。集合 X 中最小的等价关系是 X 中的相等关系；X 中最大的等价关系是 X$\times$X。

集合中的等价关系之间存在密切的联系

X 以及 X 的某些子集的集合（称为划分）。X 的一个划分是不相交的非空子集族 e，其并集为 X。如果 R 是 X 中的等价关系，且 x 在 X 中，则 x 关于 R 的等价类是由所有使得 xRy 成立的 y 所构成的集合。（由于传统习惯，这里使用“类”一词是不可避免的。）例子：若 R 是 X 中的相等关系，则每个等价类是单元素集；若 R = X $\times$ X，则 X 本身是唯一的等价类。对于 x 关于 R 的等价类，没有标准记号；我们通常将其记作 x/R，并将所有等价类的集合记作 X/R。（读作“X 模 R”，或简称为“X mod R”。习题：通过给出一个能精确指定幂集 P$($X$)$ 的子集 X/R 的条件，证明 X/R 确实是一个集合。）现在暂时忘掉 R，重新从 X 的一个划分 e 开始。在 X 中定义一个关系，我们称之为 X/e，定义为

x X/e y

当且仅当 x 和 y 属于集合族 e 的同一个集合时。我们将称 X/e 为由划分 e 导出的关系。

在上一段中，我们看到了如何为X中的每个等价关系关联一个X的子集的集合，以及如何为X的每个划分关联X中的一个关系。等价关系与划分之间的联系可以这样描述：从e到X/e的过程恰好是从R到X/R的过程的逆过程。更明确地说：如果R是X中的一个等价关系，那么等价类的集合就是X的一个划分，它诱导出关系R；如果e是X的一个划分，那么诱导出的关系是一个等价关系，其等价类的集合恰好是e。

在证明中，我们从等价关系 R 开始。由于每个 x 属于某个等价类（例如 xex/R），等价类的并集显然就是整个 X。如果 zex/R$\cap$y/R，那么 xR z 且 zRy，因此 xRy。这意味着如果两个等价类有一个公共元素，那么它们完全相同，或者换句话说，两个不同的等价类总是不相交的。因此，等价类的集合是一个划分。说两个元素属于这个划分的同一个集合（等价类），根据定义

关系，即它们彼此之间具有关系 R。这证明了我们断言的前半部分。

后一半更容易。从一个划分 e 开始, 考虑其诱导关系。由于 X 的每个元素都属于 e 的某个集合, 自反性就是说 x 和 x 在 e 的同一个集合中。对称性是说, 如果 x 和 y 在 e 的同一个集合中, 那么 y 和 x 就在 e 的同一个集合中, 这显然成立。传递性是说, 如果 x 和 y 在 e 的同一个集合中, 并且 y 和 z 在 e 的同一个集合中, 那么 x 和 z 就在 e 的同一个集合中, 这也很显然。X 中每个 x 的等价类恰好就是 x 所属的那个 e 的集合。这就完成了所承诺的全部证明。


\chapter{函数}

如果 X 和 Y 是集合，则从 X 到 Y（或映入 Y）的一个函数是一个关系 f，满足 dom f = X，并且对于 X 中的每个 x，存在 Y 中唯一的元素 y 使得 (x,y$)\in$ f。唯一性条件可以明确表述如下：若 (x,y$)\in$ f 且 (x,z$)\in$ f，则 y = z。对于 X 中的每个 x，Y 中使得 (x,y$)\in$ f 的唯一 y 记作 f$($x$)$。对于函数，这种记法及其细微变体取代了用于更一般关系的其他记法；从现在起，如果 f 是一个函数，我们将写 f$($x$)$ = y 而不再用 (x,y$)\in$ f 或 x f y。元素 y 称为函数 f 在自变量 x 处所假定（或取）的值；等价地，我们可以说 f 将 x 送到、映到或变换到 y。单词 map（或 mapping）、transformation、correspondence 和 operator 是许多有时用作 function 的同义词的词汇中的一部分。符号

f:X$\rightarrow$Y

有时被用作“f是从X到Y的函数”的缩写。所有从X到Y的函数的集合是$($X$\times$Y$)$的幂集的子集；

它将用Y表示。

上述同义词所暗示的活动性内涵使一些学者对那种认为函数无所作为、仅仅是存在的定义感到不满。这种不满反映在词汇的不同用法上：函数被保留用于指代那个以某种方式活跃的未定义对象，而我们曾称之为函数的那个有序对集合则被称作函数的图像。在数学和日常生活的精确集合论意义上，很容易找到函数的例子；我们只需寻找以表格形式呈现的信息，不一定是数值的。一个例子

是一个城市目录；在这种情况下，函数的自变量是城市的居民，而函数的值是他们的地址。

对于一般的关系，特别是函数，我们已经定义了定义域和值域的概念。从 X 到 Y 的函数 f，其定义域按定义等于 X，但它的值域不一定等于 Y；值域由 Y 中那些使得存在 X 中的 x 满足 f$($x$)$ = y 的元素 y 组成。如果 f 的值域等于 Y，我们称 f 将 X 映射到 Y 上。如果 A 是 X 的子集，我们可能想要考虑 Y 中所有那些使得存在子集 A 中的 x 满足 f$($x$)$ = y 的元素 y 的集合。这个 Y 的子集称为 A 在 f 下的像，且常记作 f$($A$)$。这个记号不好，但并非灾难性的。它不好的地方在于，如果 A 碰巧既是 X 的元素又是 X 的子集（一种不太可能但绝非不可能的情况），那么符号 f$($A$)$ 就有歧义。它是指 f 在 A 处的值，还是指 f 在 A 的元素处的值的集合？遵循通常的数学习惯，我们将使用这个不好的记号，依靠上下文，并在极少数必要时，附加口头约定以避免混淆。注意，X 本身的像就是 f 的值域；“把 X 中每个 x 映成 x”的函数 y = f$($x$)$ = x 称为 X 到 Y 的包含映射（或嵌入，或单射）。短语“由$\dots $定义的函数 f”在这种上下文中是非常常见的。当然，它的用意是暗示确实存在一个唯一的函数满足所述的条件。在眼前这个特殊情形下，这一点足够明显；我们不过是需要考虑 X$\times$Y 中所有那些使得 x = y 的有序对 (x, y$)$ 的集合。类似的考虑适用于所有情况，并且遵循通常的数学实践，我们通常通过描述函数在每个自变量 x 处的值 y 来描述一个函数。这种描述有时比直接描述所涉及（有序对）的集合更长、更繁琐，但尽管如此，大多数数学家认为“自变量-值”这种描述比其他任何描述都更加清晰。

X 到 X 的包含映射称为 X 上的恒等映射。（用关系的语言来说，X 上的恒等映射等同于 X 中的相等关系。）如果像之前那样 X $\subset$ Y，那么 X 到 Y 的包含映射与 Y 上的恒等映射之间存在一种联系；这种联系是一种从大函数构造小函数的一般步骤的特例。如果 f 是从 Y 到 Z 的函数，且 X 是 Y 的子集，那么有一种自然的方法构造从 X 到 Z 的函数 g；对于 X 中的每个 x，定义 g$($x$)$ 等于 f$($x$)$。函数 g 被称为

Y 的一个子集的包含映射是 Y 上的恒等映射在该子集上的限制。

这里是一个简单但有用的函数例子。考虑任意两个集合 X 和 Y，并通过写出 f$($x,y$)$=x 定义一个从 X$\times$Y 到 X 上的函数 f。（纯粹主义者会注意到我们应该写作 f$($(x,y$)$) 而不是 f$($x,y$)$，但没有人这样做。）函数 f 被称为从 X$\times$Y 到 X 的投影；类似地，如果 g$($x,y$)$=y，那么 g 是从 X$\times$Y 到 Y 的投影。这里的术语与之前的术语不一致，但不太严重。如果 R=X$\times$Y，那么先前称为 R 到第一坐标的投影，在现在的语言中，就是投影 f 的值域。

一个更复杂因此也更有价值的函数的例子可以如下获得。假设 R 是 X 中的一个等价关系，并设 f 是从 X 到 X/R 上的函数，定义为 f(x)=x/R。函数 f 有时被称为从 X 到 X/R 的典范映射。

如果 f 是从 X 到 Y 上的任意函数，那么存在一种自然的方式在 X 中定义一个等价关系 R；当 f(a)=f(b) 时，记作 aRb（其中 a 和 b 属于 X）。对于 Y 中的每个元素 y，令 g(y) 为 X 中所有满足 f(x)=y 的元素 x 的集合。R 的定义意味着，对于每个 y，g(y) 是关系 R 的一个等价类；换句话说，g 是从 Y 到 R 的所有等价类的集合 X/R 上的一个函数。函数 g 具有以下特殊性质：如果 u 和 v 是 Y 中不同的元素，那么 g(u) 和 g(v) 是 X/R 中不同的元素。总是将不同元素映射到不同元素的函数称为一一映射（通常称为一一对应）。在上述例子中，包含映射是一一的，但除了某些平凡的特殊情形外，投影映射不是。（习题：哪些特殊情形？）

为了介绍函数基本理论的下一个方面，我们必须暂时岔开话题，先给出最终自然数定义中的一小部分。我们认为现在没有必要定义所有自然数；我们只需要前三个。由于这不是进行冗长启发式预备知识的合适时机，我们将直接给出定义，即使可能暂时令人震惊。

或者让一些读者担忧。如下：我们通过写下来定义0、1和2

0=ø，1={x}，且2={x，{x}}。

换句话说，0是空集，1是单元素集{0}，2是对集{0,1}。

观察到在这种表面的疯狂背后自有其条理；集合0、1或2（就这些词的日常意义而言）中的元素个数分别是零、一或二。

如果 A 是集合 X 的子集，A 的特征函数是从 X 到 2 的函数 x，使得 x$($x$)$=1 或 0 分别当 x eA 或 xeX-A 时。A 的特征函数对集合 A 的依赖性可以通过写作 x\_A 而不是 x 来表示。将 X 的每个子集 A（即 o$($X$)$ 的每个元素）对应到 A 的特征函数（即 2$\times$ 的一个元素）的函数是 0$($X$)$ 与 2$\times$ 之间的一一对应。（附带说明：通常不使用“将 o$($X$)$ 中的每个 A 对应到 2$\times$ 中的元素 x\_A 的函数”这种表述，而使用缩写“函数 A $\rightarrow$ x\_A”。在这种语言中，例如，从 X$\times$Y 到 X 的投影可以称为函数 (x,y$)\rightarrow$ x，而从具有关系 R 的集合 X 到 X/R 的典范映射可以称为函数 x $\rightarrow$ x/R。）

习题。$($i$)$ Y 恰好有一个元素，即 ø，无论 Y 是否为空，并且 (ii$)$ 如果 X 非空，则 $\times$ 为空。


\chapter{族}

在某些情况下，函数的值域被认为比函数本身更重要。在这种情况下，术语和记号都会发生彻底的改变。例如，假设 z 是一个从集合 I 到集合 X 的函数。（选择这些字母本身就表明有些不同寻常的事要发生了。）定义域 I 中的一个元素称为一个指标，I 称为指标集，函数的值域称为一个加标集，函数本身称为一个族，并且函数 z 在指标 i 处的值，称为族的一项，记为 x\_i。（这些术语并非绝对固定，但在相关的细小变体中是标准选择之一；在下文中，将只使用这套术语。）一种虽不严格但普遍接受的传达記号并表明重点的方式是，谈论 X 中的一个族 $\{$z\_i$\}$，或者谈论一个由 X 的元素（无论是什么）构成的族 $\{$x$\}$；必要时，用诸如 (i$\in$I$)$ 这样的括注来表示指标集 I。因此，例如，短语“一个由 X 的子集构成的族 $\{$A\_i$\}$”通常被理解为指一个函数 A，从某个指标集 I 到 o$($X$)$。

[TRANSLATION FAILED] If{A}is a family of subsets of X,the union of the range of the family is called the union of the family {A, or the union of the sets A;;the standard notation for it is

Uiar Az 或 U:Ai,

根据是否强调指标集I而定。由并集的定义可直接得出：z$\in$UᵢAᵢ当且仅当z至少属于某一个Aᵢ。若I=2，即族$\{$Aᵢ$\}$的值域为无序对$\{$A$_0$, A$_1\}$，则UᵢAᵢ = A$_0\cup$ A$_1$。需注意，用集合族替代任意集合族来考虑并不会失去一般性；每个集合族本身就是某个指标集的值域。

对于某个族。实际上，若 e 是一个集合族，则令 e 本身充当索引集，并考虑 e 上的恒等映射作为该族。

并运算对偶所满足的代数律可以推广到任意并集。例如，设{I;}是以J为定义域的一个集族，记K=U;Ij，并令{A}是以K为定义域的一个集族。那么不难证明

这是并集结合律的推广形式。练习：阐述并证明交换律的推广形式。

空并集是有意义的（且为空集），但空交集则没有意义。除了这一细微之处，交集的术语和符号表示在所有方面都与并集完全对应。例如，若{A\_i}是一个非空集族，则该族值域的交集称为集族{A\_i}的交集，或简称为交集。

集合的标准表示法为

N:ar Ai 或 $\cap$:Ai

根据是否强调指标集I的重要性而定。（所谓“非空族”是指定义域I非空的族。）由交集的定义立即得出：若I$\neq\varnothing$，则x属于$\cap$\_$\{$i$\in$I$\}$A\_i的充要条件是x属于所有A\_i。

交集的广义交换律和结合律可以像并集那样以相同方式表述和证明，或者也可以利用德摩根定律从并集的性质推导得出。这几乎是显而易见的，因此没有太多研究价值。真正有趣的代数恒等式是那些同时涉及并集和交集的。例如，若$\{$A\_i$\}$是X的一族子集且B$\subseteq$X，则

B$\cap$U:A=U:$($B$\cap$A;$)$

和

BU$\cap$:Ai=$\cap$;$($B U A;$)$;

这些方程是分配律的一种温和推广。

习题。若{A;}与{B;}均为集合族，则

(U:A;$)\cap($U;B;$)$=Ui;$($A;$\cap$B;$)$

和

($\cap$;A;$)\cup(\cap$;B;$)$=$\cap$i,;$($A;$\cup$B;$)$

B;)。

符号说明：形如 U:,s 的记号是 U$($G,$)$eI$\times$J 的缩写形式。

族的记号是通常用于推广笛卡尔积概念时所采用的。两个集合X和Y的笛卡尔积被定义为所有有序对$($x,y$)$的集合，其中x$\in$X且y$\in$Y。这个集合与某个特定的族集合之间存在自然的——对应关系。事实上，考虑任意一个特定的无序对$\{$a,b$\}$，其中a$\neq$b，并考虑所有以$\{$a,b$\}$为指标集的族z构成的集合Z，使得zₐ$\in$X且z\_b$\in$Y。若从Z到X$\times$Y的函数f定义为f$($z$)$=$($zₐ,z\_b$)$，则f就是所承诺的一一对应。Z与X$\times$Y之间的差异仅仅是记号上的不同。笛卡尔积的推广所推广的是Z而非X$\times$Y本身。（因此，从特例过渡到一般情况时会产生一些术语上的摩擦。这是无法避免的；这正是当今数学语言的实际使用方式。）推广过程现在变得直接明了。如果$\{$Xᵢ$\}$是一个集合族（i$\in$I），那么该族的笛卡尔积按定义就是所有族$\{$xᵢ$\}$的集合，其中对于每个i$\in$I，有xᵢ$\in$Xᵢ。笛卡尔积有多种表示符号。

在目前的使用中，笛卡尔积通常用此符号表示；在本书中，我们将用

谢 X；或 X:X。

显然，若每个 \(X\_i\) 都等于同一个集合 \(X\)，则 \(\prod\_{i \in I} X\_i = X^I\)。若 \(I\) 是二元集 \(\{a,b\}\) 且 \(a \neq b\)，则通常将 \(\prod\_{i \in I} X\_i\) 等同于先前定义的笛卡尔积 \(X\_a \times X\_b\)；若 \(I\) 是单元素集 \(\{a\}\)，则类似地将 \(\prod\_{i \in I} X\_i\) 等同于 \(X\_a\) 本身。有序三元组、有序四元组等可定义为以无序三元组、四元组等为指标集的族。

设 $\{$X\_i$\}$ 是一个集合族（i$\in$I），并设 X 为其笛卡尔积。若 J 是 I 的一个子集，则对于 X 中的每个元素，自然对应着部分笛卡尔积 $\prod$\_$\{$i$\in$J$\}$ X\_i 中的一个元素。为定义此对应关系，回忆 X 中的每个元素 x 本身是一个族 $\{$x\_i$\}$，即归根结底是定义在 I 上的函数；对应的元素（记为 y）属于 $\prod$\_$\{$i$\in$J$\}$ X\_i，只需将该函数限制在 J 上即可得到。具体地，当 i$\in$J 时，记 y\_i = x\_i。对应关系 x $\rightarrow$ y 称为从 X 到 $\prod$\_$\{$i$\in$J$\}$ X\_i 的投影；我们暂时将其记为 f\_J。特别地，若 J 是单点集，例如 J=$\{$j$\}$，则我们将 f\_j（而非 f\_$\{$J$\}$）记为 f\_J。“投影”一词有多种用法；若 x$\in$X，f\_j 在 x 处的值，即 x\_j，也称为 x 在 X\_j 上的投影，或称为 x 的第 j 个坐标。定义在笛卡尔

笛卡尔积如X上的函数称为多元函数，特别地，定义在笛卡尔积Xa$\times$X$_6$上的函数称为二元函数。

练习。证明 (U\_i A\_i$)\times$ (U\_j B\_j$)$ = U\_$\{$i,j$\}$ (A\_i $\times$ B\_j$)$，并且

同样的等式对交集也成立（前提是所涉及的集族定义域非空）。同时证明（需对空集族作适当说明）：对于每个指标 \( j \)，有 \(\bigcap\_{i} X\_i \subseteq X\_j \subseteq \bigcup\_{i} X\_i\)，并且交与并实际上可被刻画为这些包含关系的极端解。这意味着：若对每个指标 \( j \) 有 \( X\_j \subseteq Y \)，则 \(\bigcup\_{i} X\_i \subseteq Y\)，且 \(\bigcup\_{i} X\_i\) 是唯一满足此极小性条件的集合；交集的表述类似。


\chapter{逆与复合}

对于任意从X到Y的函数f，存在一个从ø$($X$)$到ø$($Y$)$的函数（通常也记为f），该函数将X的每个子集A映射到Y中的像子集f$($A$)$。映射A$\rightarrow$f$($A$)$的代数性质尚不理想。若$\{$Ai$\}$是X的一族子集，则有（证明？），但相应的交集等式一般不成立（反例？），且像与补集之间的关系同样不尽如人意。

X的元素与Y的元素之间的对应关系，总是会诱导出X的子集与Y的子集之间一种性质良好的对应关系——不是通过取像的前向方式，而是通过取原像的后向方式。给定一个从X到Y的函数f，令f⁻¹（即f的逆）为从o$($Y$)$到o$($X$)$的函数，使得若B$\subseteq$Y，则$\dots $

f⁻¹$($B$)$ = $\{$ x $\in$ X : f$($x$)\in$ B $\}$。

用文字表述：\( f^{-1}(B) \) 恰好由 \( X \) 中那些被 \( f \) 映射到 \( B \) 的元素组成；集合 \( f^{-1}(B) \) 称为 \( B \) 在 \( f \) 下的原像。\( f \) 将 \( X \) 映射到 \( Y \) 上的充要条件是：\( Y \) 的每个非空子集在 \( f \) 下的原像都是 \( X \) 的非空子集。（证明？）\( f \) 为一一映射的充要条件是：\( f \) 值域中的每个单点集在 \( f \) 下的原像都是 \( X \) 中的单点集。

如果最后一个条件满足，则符号 f⁻¹ 常被赋予第二种解释，即作为一个函数，其定义域是 f 的值域，且对于 f 值域中的每个 y，其取值为 X 中满足 f(x)=y 的唯一 x。换言之，对于一一映射的函数 f，

可以写作 f⁻¹(y)=x 当且仅当。这种记法的使用是

与我们最初对f⁻¹的解释略有出入，但这种双重含义不太可能导致混淆。

像与原像之间的联系值得稍加思考。

如果BCY，则

f$($f⁻¹$($B$)$) $\subseteq$ B。

证明。若 \( y \in f(f^{-1}(B)) \)，则存在 \( x \in f^{-1}(B) \) 使得 \( y = f(x) \)；这意味着

y=f$($x$)$且f$($x$)\in$B，因此y$\in$B。

如果 f 将 X 映射到 Y 上，则

f(f⁻¹(B)) = B。

证明。若 \( y \in B \)，则存在 \( x \in X \) 使得 \( y = f(x) \)，从而存在 \( x \in f^{-1}(B) \)；这意味着 \( y \in f(f^{-1}(B)) \)。

如果选择公理成立，则

Acf⁻¹(f(A))。

证明。若 \( x \in A \)，则 \( f(x) \in f(A) \)；这意味着 \( x \in f^{-1}(f(A)) \)。

如果 \( f \) 是单射，则

A = f⁻¹(f(A))。

证明。若 \( x \in f^{-1}(f(A)) \)，则 \( f(x) \in f(A) \)，因此存在 \( u \in A \) 使得 \( f(x) = f(u) \)；由此可得 \( x = u \)，从而 \( x \in A \)。

f¹的代数性质是无可挑剔的。若$\{$B$_1\}$是Y的一族子集，则

f⁻¹$(\cup$Bᵢ$)$ = $\cup$f⁻¹$($Bᵢ$)$

和

f⁻¹$(\cap$ᵢBᵢ$)$ = $\cap$ᵢf⁻¹$($Bᵢ$)$

证明过程是直接的。例如，若 \( x \in f^{-1}(\bigcap B\_i) \)，则对所有 \( i \) 有 \( f(x) \in B\_i \)，从而对所有 \( i \) 有 \( x \in f^{-1}(B\_i) \)，因此 \( x \in \bigcap f^{-1}(B\_i) \)。该论证中的所有步骤都是可逆的。逆像的构造也与补集运算交换；即，

f⁻¹(Y-B)=X-f⁻¹(B)

对于Y的每一个子集B。实际上：若x $\in$ f⁻¹$($Y-B$)$，则f$($x$)\in$ Y-B，从而x $\notin$ f⁻¹$($B$)$，因此x $\in$ X - f⁻¹$($B$)$；这些步骤是可逆的。（注意最后一个等式实际上是一种交换律：它表明取补集后再取逆像与取逆像后再取补集是相同的。）

对逆函数的讨论表明，一个函数的作用在某种

这层含义可以保留；接下来我们将看到，两个函数的作用有时可以一步完成。具体来说，若f是从X到Y的函数，g是从Y到Z的函数，则f值域中的每个元素都属于g的定义域，因此对于X中的每个x，g$($f$($x$)$)都有意义。由h$($x$)$=g$($f$($x$)$)定义的从X到Z的函数h称为函数f与g的复合，记作g$\circ$f，或更简单地记作gf。（由于本书中无需考虑函数间的其他乘法，我们将仅使用后一种更简洁的记法。）

注意在函数复合理论中，事件的顺序至关重要。为使复合函数gf有定义，f的值域必须包含于g的定义域，但这并不必然同时保证反向成立。即使fg和gf均有定义（例如当f将X映射到Y，g将Y映射到X时），函数fg与gf也未必相同；换言之，函数复合运算不一定满足交换律。

函数复合不一定满足交换律，但总是满足结合律。若 f 将 X 映射到 Y，g 将 Y 映射到 Z，h 将 Z 映射到 U，则可构造 h 与 gf 的复合以及 hg 与 f 的复合；通过简单验证可知，两种情形得到的结果相同。

逆与复合之间的联系非常重要；类似的情形在数学中随处可见。若 f 将 X 映射到 Y，g 将 Y 映射到 Z，则 f⁻¹ 将 o$($Y$)$ 映射到 o$($X$)$，g⁻¹ 将 o$($Z$)$ 映射到 o$($Y$)$。在此情形下，可形成的复合为 gf 与 f⁻¹g⁻¹；断言后者是前者的逆。证明：若 x $\in$ (gf$)$⁻¹$($C$)$，其中 x $\in$ X 且 C $\subseteq$ Z，则 g$($f$($x$)$) $\in$ C，从而 f$($x$)\in$ g⁻¹$($C$)$，因此 x $\in$ f⁻¹$($g⁻¹$($C$)$)；论证的每一步均可逆。

函数的逆与复合是关系上类似运算的特例。因此，特别地，对于从X到Y的每个关系R，存在从Y到X的逆（或逆序）关系R⁻¹；根据定义，y R⁻¹ x 意味着 x R y。例如：若R是从X到o$($X$)$的属于关系，则R⁻¹是从o$($X$)$到X的包含关系。由相关定义可直接推出：dom R⁻¹ = ran R 且 ran R⁻¹ = dom R。若关系R是一个函数，则等价表述x R y与y R⁻¹ x可分别写作R$($x$)$=y与x$\in$R⁻¹$(\{$y$\})$。

由于交换性带来的困难，函数复合的推广需要谨慎处理。关系R与S的复合定义为：当R是从X到Y的关系，且S是从Y到Z的关系时，其复合关系记为S$\circ$R。

从Y到Z的关系。复合关系T从X到Z，记作S$\cdot$R，或简记为SR；其定义为：x T z当且仅当存在Y中的一个元素y，使得xRy且ySz。举一个具有启发性的例子，设R表示“儿子”，S表示“兄弟”，均定义在人类男性集合中。换言之，xRy表示x是y的儿子，ySz表示y是z的兄弟。此时复合关系SR表示“侄子”。（提问：R⁻¹、S⁻¹、RS和R⁻¹S⁻¹分别表示什么？）如果R和S都是函数，则xRy和ySz可分别改写为R$($x$)$=y和S$($y$)$=z。由此可知，S$($R$($x$)$)=z当且仅当x T z，因此函数复合实际上是所谓相对积的一种特例。

关系的逆与复合的代数性质与函数相同。因此，特别地，复合仅偶然满足交换律，但始终满足结合律，并且始终通过等式 (SR)⁻¹ = R⁻¹S⁻¹ 与逆运算相关联。（证明？）

关系代数提供了一些有趣的公式。假设我们暂时只考虑单一集合X中的关系，并特别令I为X中的相等关系（即X上的恒等映射）。关系I充当乘法单位元，这意味着对于X中的每个关系R，有IR = RI = R。问题：I、RR⁻¹和R⁻¹R之间是否存在某种联系？等价关系的三个定义性质可以用代数术语表述如下：自反性意味着I $\subseteq$ R，对称性意味着R $\subseteq$ R⁻¹，传递性意味着RR $\subseteq$ R。

练习。（假设每种情况下 \( f \) 是从 \( X \) 到 \( Y \) 的函数。） (i) 若 \( g \) 是从 \( Y \) 到 \( X \) 的函数，且 \( gf \) 是 \( X \) 上的恒等映射，则 \( f \) 是单射，且 \( g \) 将 \( Y \) 映射到 \( X \) 上。 (ii) 对所有 \( X \) 的子集 \( A \) 和 \( B \)，有 \( f(A \cap B) = f(A) \cap f(B) \) 的充要条件是 \( f \) 为单射。 (iii) 对所有 \( X \) 的子集 \( A \)，有 \( f(X - A) \subseteq Y - f(A) \) 的充要条件是 \( f \) 为单射。 (iv) 对所有 \( X \) 的子集 \( A \)，有 \( Y - f(A) \subseteq f(X - A) \) 的充要条件是 \( f \) 将 \( X \) 映射到 \( Y \) 上。

“二”是多少？更一般地，我们该如何定义数？为回答这个问题，让我们考虑一个集合X，并构造所有无序对$\{$a,b$\}$的汇集P，其中a属于X，b属于X，且a$\neq$b。显然，P中的所有集合都有一个共同性质，即由两个元素构成的性质。我们很容易试图将“二性”定义为P中所有集合的共同性质，但这种诱惑必须抵制——这样的定义终究是数学上的无稽之谈。什么是“性质”？我们如何知道P中所有集合只有一个共同性质？

经过一番思考，我们或许能找到一种方法，既能保留所提定义背后的思想，又无需使用诸如“共同性质”这类模糊表述。在数学实践中，将性质等同于集合——即所有具有该性质的对象构成的集合——是极为普遍的做法；为何不在此处也采用这一做法呢？换言之，为何不将“二”定义为集合P？类似的做法有时确实存在，但并非完全令人满意。问题在于，我们当前修改后的提议依赖于P，因而最终依赖于X。充其量，该提议只能为X的子集定义“二性”；它并未提示我们，何时可以将“二性”赋予一个不属于X的集合。

有两种解决途径。其一是放弃对特定集合的限制，转而考虑所有满足 a$\neq$b 的无序对 $\{$a,b$\}$。这些无序对并不构成一个集合；若要将"二"的定义建立于其上，则需将当前理论体系扩展至包含另一理论中的"非集合"（类）。此方法虽可行，但此处不予采用；我们将另辟蹊径。

数学家会如何定义一米？类似的过程是

第十一节 数 43

上述方法涉及以下两个步骤。首先，选取一个对象，该对象是所定义概念的目标模型之一——换言之，基于直觉或实践依据，如果存在任何物体值得被称为一米长，那么这个对象就应如此。其次，构建宇宙中所有与所选对象长度相同的物体的集合（注意，这并不依赖于对"米"这一概念的理解），并将"一米"定义为如此形成的集合。

一米究竟是如何定义的？选择这个例子的目的在于，对该问题的回答应能启发我们思考数字的定义方式。关键在于，在通常的一米定义中，第二步被省略了。通过某种或多或少随意的约定，选定一个物体并将其长度称为一米。如果该定义被指责为循环论证（“长度”是什么意思？），那么它可以轻易地转化为一个无可挑剔的指示性定义：毕竟没有什么能阻止我们将一米定义为等于所选定的物体。如果采用这种指示性方法，那么解释何时应将“一米性”赋予其他物体就同样简单——即当新物体与所选标准具有相同长度时。我们再次指出，判断两个物体是否具有相同长度仅依赖于简单的比较行为，并不依赖于对长度有精确的定义。

基于上述考虑，我们之前将2定义为某个特定的集合（直观上恰好有两个元素）。那么，这个标准集合是如何选定的？其他数字对应的标准集合又该如何选择？对于这个问题，数学上并没有令人信服的理由偏好某一种答案；这很大程度上取决于个人品味。选择应当以简洁和经济为原则。为了说明通常采用的特定选择，假设某个数字（比如7）已经被定义为一个集合（包含七个元素）。在这种情况下，我们该如何定义8？换句话说，我们到哪里去找一个恰好由八个元素组成的集合？我们可以在集合7中找到七个元素；那么，该用什么作为第八个元素加入其中？一个合理的答案是数字（集合）7本身。因此，建议将8定义为由7的七个元素加上7本身组成的集合。注意，根据这一建议，每个数字都将等于其所有前驱数字组成的集合。

前一段内容引出了一个对每个集合都有意义的集合论构造，但该构造仅在数的构建中具有意义。对于每个集合 \( x \)，我们定义其后继 \( x^+ \) 为将 \( x \) 本身作为元素添加到 \( x \) 的元素中所得到的集合；换言之，

（x的后继通常记作x'。）

现在我们准备定义自然数。在将0定义为一个不含任何元素的集合时，我们别无选择；必须（如我们之前所做的那样）写作

0=x.

如果每个自然数都等于其前驱集合，那么我们在定义1、2或3时也没有选择；我们必须写作

1=0+(={0})

2=1+（={0,1}），

3=2+({0,1,2})

等等。这里的“等等”表示我们采用通常的记号，并且在接下来的内容中，我们将自由使用诸如“4”或“956”这样的数字，而无需进一步解释或致歉。

根据目前所述，并不能推断出后继的构造可以在同一个集合内无限进行。我们需要一条新的集合论原理。

无穷公理。存在一个集合，它包含0，并且包含其每个元素的后继。

该公理名称的由来应当显而易见。我们尚未给出无穷的精确定义，但诸如无穷公理所描述的那类集合，似乎理应被称为无穷集。

我们暂时称集合A为后继集，如果0$\in$A，且每当x$\in$A时必有x⁺$\in$A。在此术语下，无穷公理简单地说就是存在一个后继集A。由于每个（非空）后继集族的交集本身也是后继集（请自证），因此包含在A中的所有后继集的交集构成一个后继集$\omega$。集合$\omega$是每个后继集的子集。事实上，若B是任意后继集，则A$\cap$B也是后继集。由于A$\cap$B$\subseteq$A，故A$\cap$B是定义$\omega$时所考虑的集合之一；由此可得$\omega\subseteq$A$\cap$B，从而$\omega\subseteq$B。如此建立的极小性唯一地刻画了$\omega$；外延公理保证了存在唯一一个包含于所有其他后继集的后继集。自然数定义为极小后继集$\omega$中的元素。这一定义是直观描述（即自然数由0,1,2,3,$\dots $组成）的严格对应。

第 1l 节 数 45

等等。”顺便提一下，我们用来表示所有自然数集合的符号（$\omega$）在相关著作的作者中获得了不少支持，但远未达到明确多数。在本书中，该符号将系统且唯一地用于上述定义的意义。

读者在自然数定义方面可能感到的轻微不适感相当普遍，且多数情况下是暂时的。问题在于，如同之前（在有序对的定义中）一样，被定义的对象带有一些无关的结构，这些结构看似碍事（实则无害）。我们希望被告知7的后继是8，但被告知7是8的子集或7是8的元素则令人不安。我们将利用自然数的这种上层结构，仅持续到推导出其最重要的自然性质为止；此后，这种上层结构便可安全地遗忘。

一个以自然数或全体自然数为指标集的族{x;}称为一个序列（分别为有限序列或无限序列）。若{A}是一个集合序列，其指标集为自然数n+，

则序列的并集记为

U”\_o A;or Ao U$\dots $U An.

若指标集为$\omega$，则记法为

U:-o Ai 或 AoU A$_1$U A$_2$ U $\dots $ .

序列的交集和笛卡尔积用类似的方式表示：$\cap$:\_o Ai, Ao$\cap\dots $nA,

以及

$\cap$:-o 表示 Ai, Ao $\cap$ A$_1\cap$ A$_2\cap\dots $ ，X:-o 表示 Ai, Ao $\times$ A$_1\times$ A$_2\times\dots $ 。

“序列”一词在数学文献中有几种不同的用法，但它们之间的差异更多是符号上的而非概念上的。最常见的变体是从1开始而非0；换言之，它指的是索引集为$\omega$-$\{$0$\}$而非$\omega$的族。


\chapter{皮亚诺公理}

现在我们暂时偏离主题，进行一段简短的讨论。这一讨论的目的是与自然数的算术理论进行短暂的接触。从集合论的角度来看，这算是一种令人愉悦的奢侈。

关于全体自然数构成的集合 \( \omega \)，我们所知的最重要性质是：它是唯一的、且为每个后继集合的子集的后继集合。称 \( \omega \) 为后继集合，即意味着

(I$)$ 0$\in$w

（其中，当然，0 = $\times$），且

(II) 若 new，则 n+ew

(其中 n⁺ = n $\cup\{$n$\})$。$\omega$ 的最小性性质可以表述为：如果 $\omega$ 的一个子集 S 是后继集，则 S = $\omega$。或者，用更基本的术语来说，

（III）若 SCw，且 0eS，并且每当 neS 时必有 n+eS，则 S=w。

性质$($III$)$被称为数学归纳法原理。现在我们将在此$\omega$的性质列表中添加另外两条：

(IV$)$ 对于所有 w 中的 n，n+$\neq$0。

和

(V$)$ 若 n 与 m 属于 $\omega$，且 n⁺ = m⁺，则 n = m。

(IV)的证明是平凡的；由于n+总是包含n，且由于0是空集，显然n+与0不同。(V)的证明并不平凡；它依赖于几个辅助命题。第一个命题断言某种本不应发生的情况确实不会发生。

如果证明所涉及的考虑因素看起来有些病态，且与我们在自然数理论中期望的算术精神格格不入，那么结果证明了手段的合理性。第二个命题所指的行为与刚刚排除的情况非常相似。然而，这一次，这些看似人为的考虑最终得出了肯定的结论：某种令人略感惊讶的事情总是会发生。陈述如下：（i）没有自然数是其任何元素的子集，且（ii）自然数的每个元素都是它的子集。有时，一个具有包含（$\subset$）其所有元素（$\in$）性质的集合被称为传递集。更准确地说，说E是传递的，意味着如果x$\in$y且y$\in$E，则x$\in$E。（回想一下我们在关系理论中遇到的这个词的略有不同的用法。）用这个术语来说，（ii）表明每个自然数都是传递的。

(i$)$的证明是数学归纳法原理的一个典型应用。设S为所有不属于其任何元素的自然数n的集合。（明确地说：n$\in$S当且仅当n$\in\omega$且对于任意x$\in$n，n不是x的子集。）由于0不是其任何元素的子集，因此可得0$\in$S。现在假设n$\in$S。因为n是n的子集，我们可以推断n不是n的元素，从而n⁺不是n的子集。那么n⁺可能是哪个集合的子集？如果n⁺$\subset$x，则n$\subset$x，因此（由于n$\in$S）x$\notin$n。由此可知，n⁺不能是n的子集，也不能是n的任何元素的子集。这意味着n⁺不能是n⁺的任何元素的子集，因此n⁺$\in$S。于是所需结论$($i$)$可由$($III$)$推出。

(ii$)$的证明也是归纳性的。这次设S为所有传递自然数的集合。（明确地说：n$\in$S当且仅当n$\in\omega$且对于n中的每一个x，x是n的子集。）要求0$\in$S是空真成立的。现在假设n$\in$S。若x$\in$n⁺，则要么x$\in$n，要么x=n。在第一种情形下，x$\subseteq$n（因为n$\in$S），因此x$\subseteq$n⁺；在第二种情形下，x$\subseteq$n⁺的理由更为显然。由此可知，n⁺的每个元素都是n⁺的子集，换言之，n⁺$\in$S。所需结论$($ii$)$是$($III$)$的一个推论。

现在我们准备证明$($V$)$。假设n和m是自然数，且n⁺ = m⁺。由于n $\in$ n⁺，可知n $\in$ m⁺，因此要么n $\in$ m，要么n = m。类似地，要么m $\in$ n，要么m = n。若n $\neq$ m，则必有n $\in$ m且m $\in$ n。由$($ii$)$可知n是传递的，从而推出n $\in$ n。然而n $\in$ n与$($i$)$矛盾，证明完成。

断言(I)-(V)被称为皮亚诺公理；它们曾被视为一切数学知识的源泉。

它们（连同我们已经遇到的集合论原理）可以用来定义整数、有理数、实数以及复数，并推导出它们通常的算术和分析性质。这一计划不在本书的讨论范围之内；感兴趣的读者应该不难在其他地方找到并学习相关内容。

归纳法不仅常用于证明命题，也常用于定义对象。具体而言，设f是从集合X到其自身的函数，且a是X中的一个元素。我们自然希望尝试以如下方式定义X中元素的一个无穷序列$\{$u$($n$)\}$（即从$\omega$到X的函数）：令u$($0$)$=a，u$($1$)$=f$($u$($0$)$)，u$($2$)$=f$($u$($1$)$)，依此类推。若定义者被要求解释“依此类推”的含义，他可能会借助归纳法。他或许会说，这意味着我们定义u$($0$)$为a，然后归纳地，对每个n定义u$($n+$)$为f$($u$($n$)$)。这听起来似乎合理，但作为存在性论断的论证却是不充分的。数学归纳原理确实能轻易证明至多存在一个满足所有给定条件的函数，但它并未确立此类函数的存在性。我们真正需要的是以下结论。

递归定理。若 \(a\) 是集合 \(X\) 的一个元素，且 \(f\) 是从 \(X\) 到 \(X\) 的函数，则存在一个从 \(\omega\) 到 \(X\) 的函数 \(u\)，使得 \(u(0) = a\)，并且对所有 \(n \in \omega\) 有 \(u(n^+) = f(u(n))\)。

证明。回顾一下，从$\omega$到X的函数是$\omega\times$X的某种子集；我们将明确地构造u作为有序对的集合。为此，考虑所有满足以下条件的$\omega\times$X的子集A的集合e：$($0,a$)\in$A，且若$($n,x$)\in$A则$($n⁺,f$($x$)$)$\in$A。由于$\omega\times$X本身具有这些性质，集合e非空。因此，我们可以取e中所有集合的交集u。容易验证u本身属于e，接下来只需证明u是一个函数。换言之，我们要证明：对每个自然数n，至多存在一个X中的元素x使得$($n,x$)\in$u（明确地说：若$($n,x$)$和$($n,y$)$都属于u，则x=y）。证明采用归纳法。令S为所有满足以下条件的自然数n的集合：确实至多有一个x使得$($n,x$)\in$u。我们将证明0$\in$S，且若n$\in$S，则n⁺$\in$S。

0 是否属于 S？如果不是，则存在某个与 a 不同的 b，使得 (0,b$)\in$ u。考虑在这种情况下，集合 u—$\{($0,b$)\}$。注意到这个缩减后的集合仍然包含 (0,a$)$（因为 a$\neq$b），并且如果该缩减集合包含 (n,x$)$，那么它也包含 (n⁺, f$($x$)$)。第二个断言成立的原因是，由于 n⁺$\neq$0，被丢弃的元素不等于

现在假设 n $\in$ S；这意味着存在唯一的元素 z $\in$ X 使得 (n, x$)\in$ u。由于 (n, x$)\in$ u，可得 (n⁺, f$($x$)$) $\in$ u。若 n⁺ 不属于 S，则存在某个不同于 f$($x$)$ 的 y 使得 (n⁺, y$)\in$ u。考虑此时集合 u - $\{($n⁺, y$)\}$。注意到这个缩减后的集合包含 (0, a$)$（因为 n⁺ $\neq$ 0），并且若该缩减集合包含 (m, t$)$，则它也包含 (m⁺, f$($t$)$)。事实上，若 m = n，则 t 必为 x，而缩减集合包含 (m⁺, f$($t$)$) 的原因在于 f$($x$)\neq$ y；另一方面，若 m $\neq$ n，则缩减集合包含 (m⁺, f$($t$)$) 的原因在于 m⁺ $\neq$ n⁺。换言之，u - $\{($n⁺, y$)\}\in$ e。这再次与 u 是 e 中最小集合的事实矛盾，因此我们可以得出结论：n⁺ $\in$ S。

递归定理的证明已完成。递归定理的一个应用称为归纳定义。

习题。证明：若 \( n \) 是自然数，则 \( n \neq n^+ \)；若 \( n \neq 0 \)，则存在自然数 \( m \) 使得 \( n = m^+ \)。证明 \( \omega \) 是传递的。证明：若 \( E \) 是某个自然数的非空子集，则存在 \( E \) 中的元素 \( k \)，使得当 \( m \) 是 \( E \) 中不同于 \( k \) 的元素时，有 \( k \in m \)。


\chapter{算术}

自然数加法的引入是归纳定义的一个典型例子。实际上，根据递归定理，对于每个自然数 m，存在一个从 $\omega$ 到 $\omega$ 的函数 sm，使得 sm$($0$)$ = m，并且对于每个自然数 n，有 sm$($n+$)$ = (sm$($n$)$)+；根据定义，sm$($n$)$ 的值即为和 m+n。加法的一般算术性质通过反复应用数学归纳法原理来证明。例如，加法满足结合律。这意味着

(k+m)+n=k+(m+n)

对于任意自然数 \(k\)、\(m\) 和 \(n\)，证明通过对 \(n\) 进行归纳如下：由于 \((k + m) + 0 = k + m\) 且 \(k + (m + 0) = k + m\)，

entikmirun+if n(=b0y.Ifdethefinitieqonui+s tmen) n,则根据(ke+imn-)

归纳假设） = k + (m + n) + （再次根据加法定义） = k + (m + n⁺)（同上），论证至此完成。加法交换律（即对所有 m 和 n 有 m + n = n + m）的证明稍显复杂；直接尝试可能失败。技巧在于：先通过对 n 的归纳证明 (i) 0 + n = n 和 (ii) m⁺ + n = (m + n)⁺，然后借助 (i) 和 (ii) 通过对 m 的归纳证明所需的交换律等式。

类似的方法也应用于乘积和幂的定义，以及它们基本算术性质的推导中。为了定义乘法，应用递归定理构造函数 \( p\_m \)，使得 \( p\_m(0) = 0 \)，并且对于每个自然数 \( n \)，有 \( p\_m(n^+) = p\_m(n) + m \)；那么，根据定义，\( p\_m(n) \) 的值就是乘积 \( m \cdot n \)（点号常被省略）。乘法满足结合律和交换律；

这些证明是加法情形证明的直接改编。分配律（即断言：当k、m、n均为自然数时，k$\cdot($m+n$)$=k$\cdot$m+k$\cdot$n）是数学归纳法原理的另一个简单推论（对n进行归纳）。任何以这种方式完成对加法和乘法运算的推导的人，在处理指数运算时都不会遇到困难。递归定理给出了函数eₘ，使得eₘ$($0$)$=1，且对每个自然数n有eₘ$($n⁺$)$=eₘ$($n$)\cdot$m；根据定义，eₘ$($n$)$的值即为幂mⁿ。关于幂的性质的发现与建立，以及关于乘法的详细证明，均可作为练习留给读者完成。

接下来值得关注的主题是自然数集合中的序理论。为此，我们将仔细考察自然数之间的归属关系问题。形式上，我们说两个自然数 m 和 n 是可比较的，如果 m $\in$ n，或 m = n，或 n $\in$ m。断言：两个自然数总是可比较的。该断言的证明包含若干步骤；为了方便，我们先引入一些记号。对于每个 n $\in\omega$，记 S$($n$)$ 为所有与 n 可比较的 m $\in\omega$ 的集合，并令 S 为所有满足 S$($n$)$ = $\omega$ 的 n 的集合。用这些术语表述，断言即为 S = $\omega$。我们首先证明 S$($0$)$ = $\omega$（即 0 $\in$ S）。显然 S$($0$)$ 包含 0。若 m $\in$ S$($0$)$，则由于 m $\in$ 0 不可能，要么 m = 0（此时 0 $\in$ m⁺），要么 0 $\in$ m（此时同样有 0 $\in$ m⁺）。因此，在所有情况下，若 m $\in$ S$($0$)$，则 m⁺ $\in$ S$($0$)$；这证明了 S$($0$)$ = $\omega$。接下来我们通过证明若 S$($n$)$ = $\omega$，则 S$($n⁺$)$ = $\omega$ 来完成证明。0 $\in$ S$($n⁺$)$ 是显然的（因为 n⁺ $\in$ S$($0$)$）；还需证明若 m $\in$ S$($n⁺$)$，则 m⁺ $\in$ S$($n⁺$)$。由于 m $\in$ S$($n⁺$)$，因此要么 n⁺ $\in$ m（此时 n⁺ $\in$ m⁺），要么 n⁺ = m（同样成立），要么 m $\in$ n⁺。在后一种情况下，要么 m = n（此时 m⁺ = n⁺，故 m⁺ $\in$ n⁺），要么 m $\in$ n。最后一种情况又需根据 m⁺ 与 n 的关系分类：由于 m⁺ $\in$ S$($n$)$，我们必有 n $\in$ m⁺，或 n = m⁺，或 m⁺ $\in$ n。第一种可能性与当前情况（即 m $\in$ n）矛盾，因为若 n $\in$ m⁺，则要么 n $\in$ m 要么 n = m，从而无论如何 n $\subset$ m，而我们知道没有自然数是其某个元素的子集。其余两种可能性均蕴含 m⁺ $\in$ n⁺，证明完成。

前一段表明，若 m 和 n 属于 $\omega$，则三种可能性（m$\in$n, m=n, n$\in$m）中至少有一种成立；容易看出，实际上总是恰好有一种成立。（其原因是对“自然数不是其任一元素的子集”这一事实的又一次应用。）前一段的另一推论是：若 n 和 m 是不同的自然数，则

mc n 正是 n 的传递性。反之，若 mc n 且 m$\neq$n，则 n$\in$m 不可能发生（因为此时 m 将是其某一元素的子集），因此必有 m$\in$n。若 m$\in$n，或等价地，m 是 n 的真子集，则记作 m<n，并称 m 小于 n。若已知 m 小于 n 或等于 n，则记作 m$\leq$n。注意 $\leq$ 和 < 是 $\omega$ 中的关系。前者是自反的，而后者不是；两者均非对称；但都是传递的。若 m$\leq$n 且 n$\leq$m，则 m=n。

习题。证明：若 \(m < n\)，则 \(m + k < n + k\)；并证明若 \(m < n\) 且 \(k \neq 0\)，则 \(m \cdot k < n \cdot k\)。证明：若 \(E\) 是非空自然数集合，则存在 \(E\) 中的元素 \(k\)，使得对 \(E\) 中所有 \(m\) 都有 \(k \leq m\)。

两个集合E和F（不一定是$\omega$的子集）称为等势，记作E～F，如果它们之间存在一一对应。容易验证，对于某个特定集合X的子集而言，这种意义上的等势是幂集o$($X$)$上的一个等价关系。

自然数 \(n\) 的每一个真子集都等价于某个更小的自然数（即 \(n\) 中的某个元素）。这一论断的证明采用归纳法。当 \(n=0\) 时，结论平凡成立。假设该结论对 \(n\) 成立，若 \(E\) 是 \(n^+\) 的一个真子集，则要么 \(E\) 是 \(n\) 的真子集（此时归纳假设适用），要么 \(E=n\)（结论平凡成立），要么 \(n \in E\)。在最后一种情形下，在 \(n\) 中选取一个不属于 \(E\) 的数 \(k\)，并定义 \(E\) 上的函数 \(f\)：当 \(i \neq n\) 时令 \(f(i)=i\)，且 \(f(n)=k\)。显然 \(f\) 是单射，且将 \(E\) 映到 \(n\) 中。于是 \(E\) 在 \(f\) 下的像要么等于 \(n\)，要么（由归纳假设）等价于 \(n\) 中的某个元素，从而 \(E\) 本身总是等价于 \(n^+\) 中的某个元素。

一个集合可以与自身的真子集等价，这一事实令人略感震惊。例如，若定义从 $\omega$ 到 $\omega$ 的函数 f 为：对所有 $\omega$ 中的 n，令 f$($n$)$=n+1，则 f 是全体自然数集合与由非零自然数构成的真子集之间的一一对应。令人欣慰的是，尽管全体自然数集合具有这一奇特性质，但每个具体的自然数却遵循常理。换言之，若 n 为自然数，则 n 不与自身的任何真子集等价。当 n=0 时，这一结论显然成立。现在假设该结论对 n 成立，并假设 f 是从 n+ 到其真子集 E 的一一对应。若 n$\in$E，则 f 在 n 上的限制是 n 与其真子集之间的一一对应，这与归纳假设矛盾。若 n$\notin$E，则 n 与 E-$\{$n$\}$ 等价，因此根据归纳假设，

归纳假设，\$n = E - \{n\}\$。这意味着 \$E = n^+\$，这与 \$E\$ 是 \$n^+\$ 的真子集的假设相矛盾。

一个集合 \(E\) 称为有限的，如果它与某个自然数等价；否则 \(E\) 称为无限的。

练习。使用这一定义证明 $\omega$ 是无限的。

一个集合至多只能与一个自然数等势。（证明：我们知道对于任意两个不同的自然数，其中一个必是另一个的元素，因此也是其真子集；根据前一段的结论，它们不可能等势。）由此可推知，一个有限集永远不会与其真子集等势；换言之，只要局限于有限集，整体总是大于其任何部分。

练习。利用有限性定义的这一推论来证明 $\omega$ 是无限的。

由于自然数的每一个子集都等价于一个自然数，因此有限集的每一个子集也是有限的。

有限集合E中元素的个数，按定义，是唯一与E等价的自然数；我们用\#$($E$)$表示。显然，若将E与\#$($E$)$之间的对应限制于某集合X的有限子集上，则得到从幂集o$($X$)$的某个子集到w的函数。该函数与常见的集合论关系及运算有着良好的关联。例如，若E和F是有限集合且E$\subseteq$F，则\#$($E$)\leq$\#$($F$)$。（原因在于，由于E～\#$($E$)$且F～\#$($F$)$，可知\#$($E$)$与\#$($F$)$的某个子集等价。）另一个例子是：若E和F是有限集合，则E$\cup$F是有限的；此外，若E与F不相交，则\#$($E$\cup$F$)$=\#$($E$)$+\#$($F$)$。证明中的关键步骤是：若m和n是自然数，则在和m+n中m的补集与n等价；这一辅助事实的证明通过对n进行归纳完成。类似技巧可证明：若E和F是有限集合，则E$\times$F和E^F也是有限的，且\#$($E$\times$F$)$=\#$($E$)\cdot$\#$($F$)$，\#$($E^F$)$=\#$($E$)$^\#$($F$)$。

习题。有限个有限集合的并集是有限的。若 \(E\) 是有限的，则 \(\mathcal{P}(E)\) 是有限的，且进一步地，\(\\#(\mathcal{P}(E)) = 2^{\\#(E)}\)。若 \(E\) 是非空有限自然数集合，则存在 \(E\) 中的一个元素 \(k\)，使得对 \(E\) 中所有 \(m\) 均有 \(m \leq k\)。


\chapter{序}

在整个数学中，特别是将有限集合的计数过程推广到无限集合时，序理论扮演着重要角色。基本定义十分简单。唯一需要记住的是，其主要动机来源于熟悉的“小于或等于”性质，而非“小于”。这并无深奥原因，只是“小于或等于”的推广更为常见，且更易于进行代数处理。

集合X中的关系R称为反对称的，如果对于X中的任意x和y，xRy与yRx同时成立蕴含x=y。集合X中的偏序（有时简称为序）是X上自反、反对称且传递的关系。通常，在大多数集合中，对大多数偏序使用同一个符号（或其近似印刷符号）；常用符号即常见的不等号。因此，X上的偏序可定义为X中的关系$\leq$，使得对X中所有x、y、z满足：$($i$)$ x$\leq$x，$($ii$)$ 若x$\leq$y且y$\leq$x，则x=y，$($iii$)$ 若x$\leq$y且y$\leq$z，则x$\leq$z。使用“偏”这一限定词的原因是某些关于序的问题可能无法回答。若对X中任意x和y，总有x$\leq$y或y$\leq$x，则称$\leq$为全序（有时也称为简单序或线性序）。全序集常被称为链。

练习。用涉及关系R及其逆的方程来表达反称性和完全性的条件。

偏序（而非全序）最自然的例子是包含关系。

明确地：对于每个集合 \(X\)，关系 \(C\) 是幂集 \(B(X)\) 上的一个偏序；它是全序当且仅当 \(X\) 是空集或 \(X\) 是单元素集。

全序关系的一个著名例子是自然数集合上的“小于或等于”关系。一个有趣且常见的偏序关系是函数间的扩展关系。具体而言：给定集合 X 和 Y，令 F 为所有定义域包含于 X 且值域包含于 Y 的函数的集合。在 F 上定义关系 R，使得 fRg 当且仅当 dom f $\subseteq$ dom g，且对所有 x $\in$ dom f 有 f$($x$)$=g$($x$)$；换言之，fRg 表示 f 是 g 的限制，或等价地，g 是 f 的扩展。若我们回想 F 中的函数本质上是笛卡尔积 X$\times$Y 的某些子集，便可认识到 fRg 等价于 f $\subseteq$ g；扩展关系是包含关系的一个特例。

一个偏序集是一个集合及其上的一个偏序。这种“结合”的精确表述如下：一个偏序集是一个有序对\((X, \leq)\)，其中\(X\)是一个集合，而\(\leq\)是\(X\)上的一个偏序。这种定义方式在数学中非常常见；一个数学结构几乎总是一个集合“结合”某些指定的其他集合、函数和关系。使这类定义精确化的公认方法是借助有序对、三元组或任何适当的组合。但这并非唯一的方法。例如，注意到知道一个偏序就意味着知道其定义域。因此，如果我们把偏序集描述为一个有序对，就有些冗余；仅凭第二个坐标就能传达相同的信息量。然而，在语言和符号习惯上，传统总是战胜纯粹理性。公认的数学行为（通常适用于各种结构，此处以偏序集为例）是承认有序对是正确的方法，但忘记第二个坐标才是关键，而谈论时仿佛第一个坐标才是全部。遵循惯例，我们常会说“设\(X\)是一个偏序集”，而实际意思是“设\(X\)是一个偏序的定义域”。同样的语言习惯也适用于全序集，即其序关系实际上是全序的偏序集。

偏序集理论中使用的许多术语，其技术含义与日常含义非常接近，以至于几乎不言自明。具体来说，假设X是一个偏序集，且x和y是X中的元素。当x$\leq$y时，我们记y$\geq$x；换言之，$\geq$是关系$\leq$的逆关系。若x$\leq$y且x$\neq$y，则记x<y，并称x小于或小于y，或x是y的前驱。同样地，在此情形下，我们记y>x，并称y大于或大于x，或y是x的后继。关系<具有以下性质：$($i$)$ 不存在元素x和y使得x<y与y<x同时成立。

同时，$($ii$)$ 若 x<y 且 y<z，则 x<z（即 < 是传递的）。反之，若 < 是 X 中满足 (i$)$ 和 (ii$)$ 的关系，且 x$\leq$y 定义为 x<y 或 x=y，则 $\leq$ 是 X 上的一个偏序。

$\leq$ 与 < 之间的联系可以推广到任意关系。即，给定集合 X 中的任意关系 R，我们可以通过定义 xSy 当且仅当 xRy 且 x$\neq$y 来得到 X 中的关系 S；反之，给定 X 中的任意关系 S，我们可以通过定义 xRy 当且仅当 xSy 或 x=y 来得到 X 中的关系 R。为便于简洁地指代从 R 到 S 及反向的转换，我们称 S 为 R 对应的严格关系，而 R 为 S 对应的弱关系。对于集合 X 中的某个关系，若它本身是 X 上的偏序，或其对应的弱关系是偏序，则称该关系“偏序化 X”。

设 \(X\) 是一个偏序集，若 \(a \in X\)，则集合 \(\{x \in X : x < a\}\) 是由 \(a\) 确定的初始段，通常记作 \(s(a)\)。集合 \(\{x \in X : x \leq a\}\) 是由 \(a\) 确定的弱初始段，记作 \(\bar{s}(a)\)。当需要强调初始段与弱初始段的区别时，前者称为严格初始段。一般而言，“严格”与“弱”分别对应于 \(<\) 和 \(\leq\)。例如，由 \(a\) 确定的初始段可描述为 \(a\) 的所有前驱的集合，或为强调起见，称为 \(a\) 的所有严格前驱的集合；类似地，由 \(a\) 确定的弱初始段包含 \(a\) 的所有弱前驱。若 \(x \leq y\) 且 \(y \leq z\)，则称 \(y\) 位于 \(x\) 与 \(z\) 之间；若 \(x < y\) 且 \(y < z\)，则称 \(y\) 严格位于 \(x\) 与 \(z\) 之间。若 \(x < y\) 且不存在严格位于 \(x\) 与 \(y\) 之间的元素，则称 \(x\) 是 \(y\) 的紧接前驱，或 \(y\) 是 \(x\) 的紧接后继。

若 \(X\) 是一个偏序集（特别地，可以是全序集），则可能出现 \(X\) 中存在一个元素 \(a\)，使得对 \(X\) 中每个 \(x\) 都有 \(a \leq x\) 的情况。此时称 \(a\) 为 \(X\) 的最小元（最小元素、首元）。序关系的反对称性意味着：若 \(X\) 有最小元，则它唯一。类似地，若 \(X\) 中存在一个元素 \(a\)，使得对 \(X\) 中每个 \(x\) 都有 \(x \leq a\)，则 \(a\) 为 \(X\) 的最大元（最大元素、末元）；它也是唯一的（若存在）。全体自然数构成的集合 \(\omega\)（按通常的大小顺序排序）是既有首元（即 0）却无末元的偏序集实例。同一集合若采用逆序排序，则它有末元而无首元。

在偏序集中，最小元素与极小元素之间存在重要区别。若如前所述，X 是一个偏序集，则 X 中的元素 a 被称为 X 的极小元素，当且仅当 X 中不存在严格小于 a 的元素。等价地，a 是极小元素当且仅当 x $\leq$ a 蕴含

使得 \( x = a \)。例如，考虑非空集合 \( X \) 的非空子集族 \( \mathcal{E} \)，按包含关系排序。每个单点集都是 \( \mathcal{E} \) 的极小元素，但显然 \( \mathcal{E} \) 没有最小元素（除非 \( X \) 本身是单点集）。类似地，我们区分最大元素与极大元素；\( X \) 的极大元素是满足 \( X \) 中不存在严格大于 \( a \) 的元素的元素 \( a \)。等价地，若 \( a \leq x \) 蕴含 \( x = a \)，则 \( a \) 是极大的。

在一个偏序集中，元素a被称为子集E的下界，如果对于E中的每个x都有a$\leq$x；类似地，a被称为E的上界，如果对于E中的每个x都有x$\leq$a。一个集合E可能根本没有下界或上界，也可能有许多个；在后一种情况下，可能没有一个属于E本身。（举例？）设E\_是X中E的所有下界的集合，E*是X中E的所有上界的集合。刚才所说的是E\_可能是空集，或者E\_$\cap$E可能是空集。如果E\_$\cap$E非空，那么它是一个单元素集，包含E的唯一最小元素。当然，类似的论述也适用于E*。如果集合E\_包含一个最大元素a（必然唯一），那么a被称为E的最大下界或下确界。缩写g.l.b.和inf通常被使用。由于前者发音困难，甚至难以记住g.l.b.是向上（最大）还是向下（下界），我们将只使用后一种记号。因此，inf E是X中（可能不在E内）唯一的元素，它是E的一个下界，并且大于（即大于）E的所有其他下界。另一端的定义完全类似。如果E*有一个最小元素a（必然唯一），那么a被称为E的最小上界（l.u.b.）或上确界（sup）。

与偏序集相关的概念易于表述，但需要一定时间才能充分理解。建议读者自行构造大量例子，以说明偏序集及其子集可能呈现的各种性质。为辅助这一工作，我们接下来描述三个具有某些有趣性质的特殊偏序集。 (i$)$ 该集合为 $\omega\times\omega$。为避免任何可能的混淆，我们将用中性符号 R 表示即将引入的序关系。若 (a,b$)$ 和 (x,y$)$ 是自然数的有序对，则 (a,b$)$R$($x,y$)$ 的定义意味着 (2a+1$)\cdot$2 $\leq$ (2x+1$)\cdot$2。（此处不等式

ity 符号指的是自然数的常规顺序。不愿意假装对分数无知的读者会认识到，

除了符号表示外，我们刚才定义的正是通常的序关系。

(ii$)$ 集合再次为 w$\times$w。我们再次使用一个中性符号。

对于该序，记作S。若$($a,b$)$和$($x,y$)$是自然数的有序对，则$($a,b$)$S$($x,y$)$的定义意味着：要么a严格小于x（按通常意义），要么a=x且b$\leq$y。由于这种排序方式与词典中单词的排列方式相似，故称为$\omega\times\omega$的字典序。 （iii）再次考虑集合$\omega\times\omega$。当前序关系记作T，其定义为：$($a,b$)$T$($x,y$)$意味着a$\leq$x且b$\leq$y。


\chapter{选择公理}

为了获得关于偏序集的最深刻结论，我们需要一种新的集合论工具；在此，我们暂时中断序理论的进展，以便引入这一工具。

我们首先注意到，一个集合要么是空集，要么不是；如果不是空集，那么根据空集的定义，其中必存在一个元素。这一观察可以推广。若X和Y是集合，且其中一个是空集，则笛卡尔积X$\times$Y为空集。若X和Y均非空，则存在元素x$\in$X和y$\in$Y；由此可知有序对$($x,y$)$属于笛卡尔积X$\times$Y，因此X$\times$Y非空。上述讨论构成了以下断言中n=1和n=2的情形：若$\{$Xi$\}$是有限序列的集合（例如i$\in$n），则其笛卡尔积为空集的充要条件是其中至少有一个集合为空集。该断言易于通过对n进行归纳法证明（n=0的情形会涉及关于空函数的棘手论证；不感兴趣的读者可从n=1而非n=0开始归纳）。

将前一段落中非平凡部分（必要性）推广到无穷族的情形，便得到集合论中以下重要原理。

选择公理。一个非空族（每个成员均为非空集合）的笛卡尔积

空集是非空的。

换言之：若$\{$X\_i$\}$是一个由非空集合I索引的非空集族，则存在一个族$\{$x\_i$\}$，i$\in$I，使得对每个i$\in$I，有x\_i$\in$X\_i。

假设 e 是一个由非空集合构成的非空集族。我们可以将 e 视为一个族，或者更准确地说，我们可以将 e 转化为一个索引集，只需将 e 本身作为索引集，并将 e 上的恒等映射作为索引函数。该公理

选择公理进而指出，集合族 e 的笛卡尔积至少有一个元素。根据定义，这样一个笛卡尔积的元素是一个函数（族、索引集），其定义域为索引集（此处为 e），且在每个索引处的值属于该索引对应的集合。结论：存在一个定义域为 e 的函数 f，使得若 A $\in$ e，则 f$($A$)\in$ A。这一结论特别适用于 e 为非空集合 X 的所有非空子集构成的集族的情况。在该情况下，该论断

存在一个函数 f，其定义域为 o(X) 一 { }，使得若 A

若在该定义域内，则 f$($A$)\in$ A。用直观语言描述，函数 f 可视为同时从多个集合中各选取一个元素；这正是该公理名称的由来。（在此意义上，从一个集合 X 的每个非空子集中“选取”一个元素的函数，称为 X 的选择函数。）我们已知，若待选集合族为有限集，则同时选取是可能的。

是我们在选择公理提出之前就已知道的一个简单推论；该公理的作用在于保证无限情形下的这种可能性。

选择公理在前一段中的两个推论（一个针对集合的幂集，另一个针对更一般的集合族）实际上只是该公理的重新表述。过去人们认为，考察选择公理每个推论在证明中对该公理的依赖程度至关重要。若能在不借助选择公理的情况下给出替代证明，便意味着胜利；反之，若能证明该推论（在集合论其余公理的基础上）等价于选择公理，则意味着虽败犹荣。介于两者之间的情况则令人沮丧。作为示例（及练习），我们提及一个断言：每个关系都包含一个具有相同定义域的函数。另一个示例：若e是两两不相交的非空集合族，则存在一个集合A，使得对于e中的每个C，A$\cap$C恰为单元素集。这两个断言均属于已知与选择公理等价的众多命题之列。

作为选择公理应用的一个例证，考虑如下断言：如果一个集合是无限的，那么它有一个与$\omega$等价的子集。一个非正式的论证可能如下进行：如果X是无限的，那么特别地，它非空（即它不与0等价）；因此它有一个元素，记为x$_0$。由于X不与1等价，集合X-$\{$x$_0\}$非空；因此它有一个元素，记为x$_1$。无限重复这一论证；例如，下一步是说明X-$\{$x$_0$, x$_1\}$非空，因此它有一个元素，记为x$_2$。结果得到一个由X中不同元素组成的无限序列$\{$xₙ$\}$；证毕。这个证明的草图至少有一个优点，即

诚实地讲，其背后最重要的思想是：从一个非空集合中选择一个元素这一操作被无限重复进行。熟悉选择公理用法的数学家常常会提出这样非正式的论证；他的经验使他能一眼看出如何将其精确化。就我们的目的而言，更明智的做法是更仔细地审视这一过程。

设 \( f \) 为 \( X \) 的一个选择函数；即 \( f \) 是从 \( X \) 的所有非空子集构成的集族到 \( X \) 的函数，且对 \( f \) 定义域中的每个 \( A \) 有 \( f(A) \in A \)。令 \( \mathcal{E} \) 为 \( X \) 的所有有限子集构成的集族。由于 \( X \) 是无限的，因此若 \( A \in \mathcal{E} \)，则 \( X - A \) 非空，从而 \( X - A \) 属于 \( f \) 的定义域。定义函数 \( g \) 从 \( \mathcal{E} \) 到 \( \mathcal{E} \) 如下：\( g(A) = A \cup \{ f(X - A) \} \)。换言之，\( g(A) \) 是通过将 \( f \) 从 \( X - A \) 中选取的元素添加到 \( A \) 中得到的。我们对函数 \( g \) 应用递归定理；例如，我们可以从空集 \( \varnothing \) 开始启动该过程。结果存在一个从 \( \omega \) 到 \( \mathcal{E} \) 的函数 \( U \)，使得 \( U(0) = \varnothing \)，且对每个自然数 \( n \) 有 \( U(n^+) = U(n) \cup \{ f(X - U(n)) \} \)。断言：若 \( v(n) = f(X - U(n)) \)，则 \( v \) 是从 \( \omega \) 到 \( X \) 的一一对应，因此 \( \omega \) 确实与 \( X \) 的某个子集（即 \( v \) 的值域）等势。为证明该断言，我们进行一系列基本观察；其证明由定义直接可得。第一：对所有 \( n \) 有 \( v(n) \notin U(n) \)。第二：对所有 \( n \) 有 \( v(n) \in U(n^+) \)。第三：若 \( n \) 和 \( m \) 是自然数且 \( n \leq m \)，则 \( U(n) \subseteq U(m) \)。第四：若 \( n \) 和 \( m \) 是自然数且 \( n < m \)，则 \( v(n) \neq v(m) \)。（理由：\( v(n) \in U(m) \) 但 \( v(m) \notin U(m) \)。）最后的观察表明 \( v \) 将不同的自然数映射到 \( X \) 中不同的元素；我们只需记住：任意两个不同的自然数中，总有一个严格小于另一个。

证明完成；现在我们得知每一个无限集都有一个与$\omega$等价的子集。这一结果在此证明，与其说是出于其内在价值，不如说是作为选择公理正确运用的一个范例，它有一个有趣的推论。该断言是：一个集合是无限的当且仅当它与自身的某个真子集等价。“若”的部分我们已经知晓，它仅仅说明有限集不可能与真子集等价。为证明“仅当”，假设X是无限的，并设$\nu$是从$\omega$到X的一一对应。若x在$\nu$的值域中，设x=$\nu($n$)$，则记h$($x$)$=$\nu($n⁺$)$；若x不在$\nu$的值域中，则记h$($x$)$=x。容易验证h是从X到自身的一一对应。由于h的值域是X的真子集（它不包含$\nu($0$)$），该推论的证明完成。这一推论曾被戴德金用作无限的定义本身。


\chapter{佐恩引理}

存在性定理断言，在某个特定集合中存在一个具有特定性质的对象。许多存在性定理可以（或在必要时重新）表述为：其基础集合是一个偏序集，而关键性质是极大性。我们接下来的目标是陈述并证明这类定理中最重要的一个。

佐恩引理。若\$X\$是一个偏序集，且\$X\$中的每一个链

有上界，则X包含一个极大元。

讨论。回顾一下，链是全序集。所谓“X中的链”，指的是X的一个子集，且该子集本身作为偏序集时是全序的。若A是X中的一个链，佐恩引理的假设保证了A在X中存在上界，但并不保证A在A自身中存在上界。佐恩引理的结论是：X中存在一个元素a，使得若a$\leq$x，则必有a=x。

该证明的基本思路与我们之前讨论无限集时所用的方法类似。根据假设，X 非空，因此它有一个元素，记为 x$_0$。如果 x$_0$ 是极大元，则在此停止；如果不是，则存在一个严格大于 x$_0$ 的元素，记为 x$_1$。如果 x$_1$ 是极大元，则在此停止；否则继续。如此无限重复这一论证，最终必然得出一个极大元。

论证的最后一句可能是最不令人信服的部分，它隐藏了诸多困难。例如，考虑以下可能性：该论证无限重复下去，可能导致一整个无限的非极大元素序列；在这种情况下我们该如何处理？答案是，这样一个无限序列的值域是X中的一个链，因此存在一个上界；此时需要做的是，从该上界开始，将整个论证重新进行一遍。

上界。至于这一切究竟何时、以何种方式终结，至少可以说并不明确。对此别无他法，我们必须审视精确的证明。该证明的结构改编自策梅洛最初给出的一个版本。

证明。第一步是将抽象的偏序关系替换为某个合适集合族中的包含序。更精确地说，对于X中的每个元素x，我们考虑由x及其所有前驱组成的弱初始段s$($x$)$。函数s（从X到o$($X$)$）的值域s是X的某个子集族，当然我们可以将其视为按包含关系（偏）序化。函数s是一一对应的，且s$($x$)\subseteq$ s$($y$)$的充要条件是x $\leq$ y。基于此，在X中寻找极大元素的任务等价于在s中寻找极大集合。关于X中链的假设蕴含（实际上等价于）关于s中链的相应陈述。

设 \(s\_x\) 为 \(X\) 中所有链的集合；\(X\) 中的每个成员都包含于某个 \(x \in X\) 对应的 \(s(x)\) 中。该族 \(s\) 是一个非空集合族，按包含关系偏序，且满足：若 \(e\) 是 \(8c\) 中的一条链，则 \(e\) 中各集合的并集（即 \(\bigcup\_{A \in e} A\)）属于该族。由于 \(SC\) 中的每个集合都被 \(s\) 中的某个集合所支配，从 \(S\) 到该族的过渡不会引入任何新的极大元。族 \(X\) 的一个优势在于链假设具有更具体的形式：我们不再仅仅说每条链 \(e\) 在 \(S\) 中有某个上界，而是可以明确地指出，\(e\) 中各集合的并集（显然是 \(e\) 的一个上界）是族 \(9\) 中的一个元素。\(S\) 的另一个技术优势在于它包含其每个集合的所有子集；这使得我们可以逐个元素地缓慢扩大 \(X\) 中的非极大集。

现在我们可以忽略X上给定的偏序关系。接下来，我们考虑非空集合X的一个非空子集族x，它满足两个条件：x中每个集合的所有子集都属于s，并且x中每个链的并集都属于c。注意，第一个条件蕴含空集属于C。我们的任务是证明在SC中存在一个极大集。

设 \( f \) 为 \( X \) 的一个选择函数，即 \( f \) 是从 \( X \) 的所有非空子集构成的族到 \( X \) 的一个函数，且对于 \( f \) 定义域中的每个 \( A \)，有 \( f(A) \in A \)。对于 \( \mathfrak{C} \) 中的每个集合 \( A \)，令 \( A^* \) 为 \( X \) 中所有满足将 \( x \) 添加到 \( A \) 后所得集合仍属于 \( \mathfrak{C} \) 的元素 \( x \) 构成的集合；换言之，\( A^* = \{ z \in X : A \cup \{ z \} \in \mathfrak{C} \} \)。定义从 \( \mathfrak{C} \) 到 \( \mathfrak{C} \) 的函数 \( g \) 如下：若 \( A^* - A \neq \varnothing \)，则 \( g(A) = A \cup \{ f(A^* - A) \} \)；若 \( A^* - A = \varnothing \)，则 \( g(A) = A \)。由 \( A^* \) 的定义可知，\( A^* - A = X \) 当且仅当 \( A \) 是极大的。因此，用这些术语表述，我们需要证明的是：在 \( \mathfrak{C} \) 中存在一个集合 \( A \) 使得 \( g(A) = A \)。事实证明，关键的性质是：

g 的一个性质是：g(A)（总是包含 A）至多比 A 多一个元素。

现在，为便于阐述，我们引入一个临时定义。

我们称SC的一个子集族s为一个塔，如果

(i) Øe³,

(ii$)$ 若 A $\in$ 3，则 g$($A$)\in$ 5。

(iii$)$ 若 e 是 3 中的一条链，则 $\cup$Ae$\in$Ae$\in$3。

塔楼确实存在；整个集合 sx 就是一个塔楼。由于一组塔楼的交集仍然是塔楼，因此特别地，如果 Jo 是所有塔楼的交集，那么 3o 就是最小的塔楼。我们当前的目标是证明塔楼 3o 是一个链。

设称3o中的集合C为可比较的，若它与3o中的每个集合都可比较；这意味着若A$\in$3o，则要么A$\subseteq$C，要么C$\subseteq$A。称3o为链，是指3o中所有集合都是可比较的。可比较的集合必然存在；空集便是其中之一。在接下来的几段中，我们将注意力集中于一个任意但暂时固定的可比较集合C上。

假设 Ae3o 且 A 是 C 的真子集。断言：g$($A$)\subseteq$ C。理由如下：由于 C 是可比较的，要么 g$($A$)\subseteq$ C，要么 C 是 g$($A$)$ 的真子集。在后一种情况下，A 是 g$($A$)$ 的某个真子集的真子集，这与 g$($A$)$ - A 不能多于一个单元素的事实相矛盾。

接下来考虑集合 \( u \)，它由 \( \mathfrak{3}\_0 \) 中所有满足 \( A \subseteq C \) 或 \( g(C) \subseteq A \) 的集合 \( A \) 组成。集合 \( u \) 略小于 \( \mathfrak{5}\_0 \) 中与 \( g(C) \) 可比较的集合族；实际上，若 \( A \in u \)，由于 \( C \subseteq g(C) \)，则要么 \( A \subseteq g(C) \)，要么 \( g(C) \subseteq A \)。断言：\( u \) 是一个塔。由于 \( \emptyset \subseteq C \)，塔的第一个条件成立。为证明第二个条件，即若 \( A \in u \)，则 \( g(A) \in u \)，将讨论分为三种情形。第一：\( A \) 是 \( C \) 的真子集。则由前一段可知 \( g(A) \subseteq C \)，因此 \( g(A) \in u \)。第二：\( A = C \)。则 \( g(A) = g(C) \)，故 \( g(C) \subseteq g(A) \)，因此 \( g(A) \in u \)。第三：\( g(C) \subseteq A \)。则 \( g(C) \subseteq g(A) \)，因此 \( g(A) \in u \)。塔的第三个条件，即 \( u \) 中一个链的并集属于 \( u \)，由 \( u \) 的定义直接可得。结论：\( u \) 是包含在 \( \mathfrak{3}\_0 \) 中的一个塔，因此，由于 \( \mathfrak{5}\_0 \) 是最小塔，有 \( u = \mathfrak{5}\_0 \)。

前面的讨论意味着，对于每个可比较集合C，集合g$($C$)$也是可比较的。理由：给定C，按上述方式构造u；由于u=3o这一事实意味着，若A$\in$Jo，则要么A$\subseteq$C（此时A$\subseteq$g$($C$)$），要么g$($C$)\subseteq$A。

我们现在知道ø是可比较的，并且g将可比较集映射到可比较集。由于可比较集的链的并集是可比较的，因此$($在3$_0$中的$)$可比较集构成一个塔，从而它们穷尽了3$_0$；这正是我们关于3$_0$所要证明的结论。

由于3o是一个链，所有3o中集合的并集（记为A）本身是5o中的一个集合。由于该并集包含3o中的所有集合，因此有g$($A$)\subseteq$ A。又因为总有A $\subseteq$ g$($A$)$，所以可得A = g$($A$)$，佐恩引理的证明至此完成。

习题：佐恩引理等价于选择公理。[证明提示：给定一个集合X，考虑满足定义域dom f $\subseteq$ (X$)$、值域ran f $\subseteq$ X且对于所有A$\in$dom f有f$($A$)\in$A的函数f；通过扩展对这些函数进行排序，利用佐恩引理找出其中的一个极大函数，并证明若f是极大的，则dom f = (X$)$ - $\{\varnothing\}$。] 考虑以下每个陈述，并证明它们也等价于选择公理。$($i$)$ 每个偏序集都有一个极大链（即一个不是任何其他链的真子集的链）。

每个偏序集中的链都包含于某个极大链中。(iii) 每个链都有最小上界的偏序集必有一个极大元素。


\chapter{良序}

一个偏序集可能没有最小元素；即使有，也完全有可能某个子集没有最小元素。如果一个偏序集的每个非空子集都有最小元素，则称该偏序集为良序的（其序关系称为良序）。这个定义的一个推论——在查看任何例子和反例之前就值得注意——是每个良序集都是全序集。事实上，如果x和y是良序集中的元素，那么$\{$x, y$\}$是该良序集的一个非空子集，因此有第一个元素；根据第一个元素是x还是y，我们得到x $\leq$ y或y $\leq$ x。

对于每个自然数 \(n\)，其所有前驱的集合（即根据我们的定义，集合 \(n\)）是一个良序集（按大小排序），而所有自然数的集合 \(\omega\) 也是如此。集合 \(\omega \times \omega\)，其中 \((a,b) \leq (x,y)\) 定义为 \((2a+1)2 \leq (2x+1)2\)，不是良序的。一个验证方法是注意到对所有 \(a\) 和 \(b\) 有 \((a,b+1) \leq (a,b)\)；由此可知整个集合 \(\omega \times \omega\) 没有最小元素。然而，\(\omega \times \omega\) 的某些子集确实有最小元素。例如，考虑所有满足 \((1,1) \leq (a,b)\) 的序对 \((a,b)\) 构成的集合 \(E\)；集合 \(E\) 以 \((1,1)\) 为其最小元素。注意：若将 \(E\) 视为其自身的偏序集，它仍然不是良序的。问题在于，尽管 \(E\) 有最小元素，但 \(E\) 的许多子集却没有最小元素；例如，考虑 \(E\) 中所有满足 \((a,b) \neq (1,1)\) 的序对 \((a,b)\) 构成的集合。再举一例：\(\omega \times \omega\) 在其字典序下是良序的。

关于良序集的一个最令人愉快的事实是，我们可以通过类似于数学归纳法的过程来证明其元素的性质。具体来说，假设 \( S \) 是良序集 \( X \) 的一个子集，并且假设对于 \( X \) 中的任意元素 \( x \)，若整个初始段 \( s(x) \) 都包含在 \( S \) 中，则 \( x \) 本身也属于 \( S \)；超限归纳原理断言，在此条件下，我们必然有

设 S=X。等价地：如果在一个集合中，某个元素的所有严格前驱都存在必然意味着该元素本身也存在，那么这个集合必定包含一切。

在审视证明之前，有几点需要说明。普通数学归纳法原理与超限归纳法原理在表述上有两个显著区别。其一：后者并非从前驱元素过渡到每个元素，而是从所有前驱元素的集合过渡到每个元素。其二：后者不假设存在起始元素（如零）。第一个区别至关重要：良序集中的某个元素可能没有直接前驱。当本表述应用于$\omega$时，可轻易证明其与数学归纳法原理等价；然而，该原理应用于任意良序集时，并不等价于超限归纳法原理。换言之，这两个表述通常互不等价；它们在$\omega$中的等价性是一种幸运但特殊的状况。

这里有一个例子。设 X 为 $\omega$+，即 X = $\omega\cup\{\omega\}$。在 X 中定义序：按通常方式排列 $\omega$ 中的元素，并要求对所有 n $\in\omega$ 有 n < $\omega$。结果得到一个良序集。问题：是否存在 X 的一个真子集 S，使得 0 $\in$ S，并且每当 n $\in$ S 时都有 n+1 $\in$ S？答案：存在，即 S = $\omega$。

普通归纳与超限归纳的第二个区别（后者不需要起始元素）更多是语言层面而非概念层面的。若 x$_0$ 是 X 的最小元素，则 s$($x$_0)$ 为空集，因此 s$($x$_0)\subseteq$ S；超限归纳原理的假设条件因而要求 x$_0$ 属于 S。

超限归纳原理的证明几乎是平凡的。若X-S非空，则它有一个最小元素，记为x。这意味着初始段s(x)中的每个元素都属于S，因此根据归纳假设，x也属于S。这就产生了矛盾（x不可能同时属于S和X-S）；结论是X-S最终为空集。

称良序集A为良序集B的延续，需满足：首先，B是A的子集；其次，B是A的初始段；最后，B中元素的序关系与A中相同。因此，若X是良序集，且a与b为X中满足b<a的元素，则s(a)是s(b)的延续，且显然X是s(a)与s(b)的延续。

如果 e 是良序集的任意一个由初始段构成的集族，则 e 在延续关系下构成一个链；这意味着 e 是一个

良序集的一个性质是：在该集合的任意两个不同成员中，其中一个必是另一个的延续。这一论断的某种逆命题同样成立且经常有用。若一族良序集 e 关于延续关系构成一个链，且 U 是 e 中所有集合的并集，则存在唯一的 U 的良序，使得 U 是 e 中每个（不同于 U 自身的）集合的延续。粗略地说，良序集链的并集是良序的。这种简略表述存在风险，因为它未说明"链"是相对于延续关系而言的。若将"链"一词隐含的序关系理解为保序包含，则结论不成立。

证明是直接的。若a和b属于U，则存在e中的集合A和B，使得a$\in$A且b$\in$B。由于要么A=B，要么A和B中一个是另一个的延续，因此在所有情况下，a和b都属于e中的某个集合；U上的序通过以下方式定义：对每一对$\{$a,b$\}$，按照e中同时包含a和b的任意集合中的顺序来排序。由于e是一个链，这个顺序是唯一确定的。（定义U中预期顺序的另一种方法是：回忆e中集合的给定顺序是有序对的集合，然后取所有这些有序对集合的并集。）

直接验证表明，上一段定义的关系确实是一个序，而且其构造在每一步都是必然的（即最终序由给定的序唯一确定）。证明该结果实际上是一个良序同样直接。U的每个非空子集必然与e中某个集合有非空交集，因此在该集合中必有首元素；而e是延续链这一事实意味着该首元素也必然是U的首元素。

习题。设偏序集 \(X\) 的子集 \(A\) 是 \(X\) 中的共尾子集，若对 \(X\) 中的每个元素 \(x\)，都存在 \(A\) 中的元素 \(a\) 使得 \(x \leq a\)。证明：每个全序集都有一个共尾的良序子集。

良序的重要性源于以下结论，由此我们可以推断出，超限归纳原理的适用范围远比表面看起来要广泛得多。

良序定理：每个集合都可以被良序化。

讨论。一个更好（但不太传统）的表述是：对于每个集合X，存在一个定义域为X的良序。注意：该良序

第 17 节 良序 69

并不保证与给定集合可能已经拥有的任何其他结构有任何关系。例如，如果读者知道某些偏序集或全序集的序关系明确不是良序，他不应草率地认为自己发现了悖论。唯一能得出的结论是：某些集合可以以多种方式排序，其中一些是良序而另一些不是——这一点我们早已知道。

证明. 我们应用佐恩引理。给定集合X，考虑X的所有良序子集构成的族W。具体而言：W的一个元素是X的子集A连同A上的一个良序。我们通过延续关系对W进行偏序。

集合族 W 非空，例如 Ø $\in$ W。若 X $\neq$ Ø，则可构造 W 中更简单的元素，例如对 X 中任意特定元素 x，$\{($x,x$)\}$ 即为一例。若 e 是 W 中的一个链，则 e 中所有集合的并集 U 具有唯一的良序，使得 U "大于"（或等于）e 中的每个集合——这正是我们之前关于延续的讨论所实现的结果。这意味着佐恩引理的主要前提已得到验证；其结论是：在 W 中存在一个极大良序集，记为 M。该集合 M 必须等于整个集合 X。理由：若 x 是 X 中不属于 M 的元素，则可将 x 置于 M 的所有元素之后来扩充 M。这一明确但非形式化描述的严格表述留作读者练习。至此，良序定理的证明完成。

习题。证明：一个全序集是良序的当且仅当每个元素的严格前驱集是良序的。对于偏序集是否也存在类似的条件？证明良序定理蕴含选择公理（因此它与该公理以及佐恩引理等价）。证明：若 \( R \) 是集合 \( X \) 上的一个偏序，则存在 \( X \) 上的一个全序 \( S \) 使得 \( R \subseteq S \)；换言之，每个偏序都可以在不扩大定义域的情况下扩展为一个全序。


\chapter{超限递归}

“归纳定义”过程有一个超限类比。

普通递归定理在$\omega$上构造一个函数；其原始材料是从$\omega$中每个非零元素n的前驱元素的值得到该函数在n处的值的方法。超限类比则在任何良序集W上构造一个函数；其原始材料是从W中每个元素a的所有前驱的值得到该函数在a处的值的方法。

为简洁表述该结果，我们引入一些辅助概念。若a是良序集W中的元素，X为任意集合，则称从W中a的初始段到X的函数为X中a型序列。当a$\in\omega$⁺时，a型序列即此前所称的序列——根据a<$\omega$或a=$\omega$分别对应有限或无限序列。若U是从W到X的函数，则对每个a$\in$W，U在a的初始段s$($a$)$上的限制即为a型序列的实例；下文将用U（而非U|s$($a$)$）表示该序列以方便论述。

类型为W在X中的序列函数是一个函数f，其定义域由X中所有类型为a的序列（a为W中的任意元素）组成，且值域包含于X。粗略地说，序列函数告诉我们如何“延长”一个序列：给定一个延伸至W中某个元素（但不包含该元素）的序列，我们可以利用序列函数在其末尾添加一项。

超限递归定理。若W是一个良序集，且f是

设 \( W \) 为 \( X \) 上的一个序列函数，则存在唯一的函数

存在一个从 W 到 X 的 U，使得对于 W 中的每个 a，有 U(a) = f(Ua)。

证明。唯一性的证明是一个简单的超限归纳法。为证明存在性，回忆从W到X的函数是W$\times$X的某种特定子集；我们将明确地将U构造为一组有序对。称W$\times$X的子集A为f-封闭的，如果它满足以下性质：每当a$\in$W且t是包含在A中的类型为a的序列（即对于s$($a$)$初始段中的所有c，有$($c, t$($c$)$)$\in$A），则$($a, f$($t$)$)$\in$A。由于W$\times$X本身是f-封闭的，这样的集合确实存在；令U为所有这些集合的交集。由于U本身是f-封闭的，只需证明U是一个函数。换言之，我们要证明：对于每个c$\in$W，至多存在一个x$\in$X使得$($c, x$)\in$U。（明确地说：如果$($c, x$)$和$($c, y$)$都属于U，则x = y。）证明采用归纳法。令S为W中所有满足以下条件的元素c的集合：确实至多存在一个x使得$($c, x$)\in$U。我们将证明：如果s$($a$)\subseteq$ S，则a$\in$S。

设 s$($a$)\subset$ S 意味着：若在 W 中 c < a，则存在唯一的 X 中元素 x 使得 (c, x$)\in$ U。由此定义的对应 c $\rightarrow$ x 是一个类型为 a 的序列，记为 t，且 t $\subset$ U。若 a 不属于 S，则存在某个不同于 f$($t$)$ 的 y 使得 (a, y$)\in$ U。断言：集合 U - $\{($a, y$)\}$ 是 f-封闭的。这意味着：若 b $\in$ W 且 r 是包含于 U - $\{($a, y$)\}$ 的类型为 b 的序列，则 (b, f$($r$)$) $\in$ U - $\{($a, y$)\}$。事实上，若 b = a，则 r 必为 t（由定理的唯一性断言），而缩减后的集合包含 (b, f$($r$)$) 的原因是 f$($t$)\neq$ y；另一方面，若 b $\neq$ a，则缩减后的集合包含 (b, f$($r$)$) 的原因是 U 是 f-封闭的（且 b $\neq$ a）。这与 U 是最小 f-封闭集的事实矛盾，因此可得出结论：a $\in$ S。

超限递归定理存在性断言的证明至此完成。超限递归定理的一个应用被称为超限归纳定义。

我们继续讨论序理论中的一个重要部分，顺便提一下，这部分内容也将说明超限递归定理如何应用。

两个偏序集（特别地，可以是全序集甚至良序集）称为相似的，如果它们之间存在一个保序的一一对应。更明确地说：称偏序集X和Y相似（记作X≌Y）意味着存在一个从X到Y上的一一对应（记为f），使得对于X中的任意元素a和b，f$($a$)\leq$f$($b$)$（在Y中）成立的充要条件是a$\leq$b（在X中）。这样的对应f有时被称为相似映射。

习题。证明相似性保持<（与定义要求保持$\leq$的意义相同），并且实际上，

一个将偏序集映射到另一个偏序集的一一对应函数是相似映射当且仅当它保持<

偏序集X上的恒等映射是一个相似映射。

X 到 X。如果 X 和 Y 是偏序集，且当且仅当 f 是一个相似映射

从X到Y的满射，那么（由于f是单射）存在一个明确的

确定从Y到X上的逆函数f⁻¹，且f⁻¹是一个相似映射。

此外，若 g 是从 Y 到偏序集 Z 上的相似映射，则

复合映射 gf 是从 X 到 Z 的一个相似映射。由此可得

注释：若我们将注意力限制在某个特定集合E上，且f，

因此，我们仅考虑定义域为子集的偏序关系。

若在集合 \(E\) 上，则相似性是偏序集集合中的一个等价关系。

这样得到的集合也是如此。即使我们进一步缩小范围，情况依然不变。

在良序集集合上如此得到的一个等价关系。我们只考虑定义域包含在相似性中的良序。

在良序集合中得到的等价关系。

尽管相似性是在完全一般性中对偏序集定义的，

一般性，并且可以在该层面上进行研究，我们对以下内容的兴趣

以下内容仅适用于良序集合的相似性。

一个良序集完全有可能与其真子集相似；例如，考虑全体自然数构成的集合和全体偶数构成的集合。（一如既往，自然数 m 被定义为偶数当且仅当存在自然数 n 使得 m=2n。映射 n$\rightarrow$2n 是从全体自然数集到全体偶数集的一个相似映射。）然而，一个良序集与其自身的一部分的相似映射是一种非常特殊的映射。事实上，若 f 是良序集 X 到自身的一个相似映射，则对 X 中的每个元素 a 都有 a$\leq$f$($a$)$。该证明直接基于良序的定义。如果存在元素 b 使得 f$($b$)$<b，那么其中必有一个最小者。设 b 为该最小者，若 a<b，则 a$\leq$f$($a$)$；特别地，取 a=f$($b$)$，可得 f$($b$)\leq$f$($f$($b$)$)。然而，由于 f$($b$)$<b，且 f 保持序关系，故有 f$($f$($b$)$)<f$($b$)$。要摆脱这一矛盾，唯一途径是承认 f$($b$)$<b 不可能成立。

上一段的结果有三个特别有用的推论。第一个是：如果两个良序集 \(X\) 和 \(Y\) 是相似的，那么它们之间恰好存在一个相似映射。假设 \(g\) 和 \(h\) 都是从 \(X\) 到 \(Y\) 的相似映射，并令 \(f = g^{-1}h\)。由于 \(f\) 是 \(X\) 到自身的相似映射，因此对 \(X\) 中的每个 \(a\) 有 \(a \leq f(a)\)。这意味着对 \(X\) 中的每个 \(a\) 有 \(a \leq g^{-1}(h(a))\)。应用 \(g\)，我们推断出对 \(X\) 中的每个 \(a\) 有 \(g(a) \leq h(a)\)。由于 \(g\) 和 \(h\) 的地位对称，我们同样可以推断出对 \(X\) 中的每个 \(a\) 有 \(h(a) \leq g(a)\)。结论：\(g = h\)。

第二个推论是：良序集永远不会与其某个初始段相似。事实上，若 \(X\) 是良序集，\(a\) 是 \(X\) 中的元素，且 \(f\) 是从 \(X\) 到 \(s(a)\) 的相似映射，则特别地有 \(f(a) \in s(a)\)，从而 \(f(a) < a\)，这是不可能的。

第三个且最重要的结论是良序集的可比较性定理。该断言指出：若X与Y均为良序集，则要么X与Y相似，要么其中一个与另一个的初始段相似。为练习起见，我们将在证明中使用超限递归定理，尽管完全可避免使用它。假设X与Y是非空良序集，且两者均不与对方的初始段相似；我们将在这种情况下证明X必与Y相似。设a$\in$X，且t是X中类型为a的序列——换言之，t是从s$($a$)$到Y的函数。令f$($t$)$为t的值域在Y中的最小真上界（若存在）；否则，令f$($t$)$为Y的最小元素。用超限递归定理的术语来说，由此确定的函数f是X到Y中的序列函数。设U为超限递归定理在此情形下关联的函数。通过简单论证（使用超限归纳法）可知，对X中每个a，函数U将X中由a确定的初始段一一映射到Y中由U$($a$)$确定的初始段上。这意味着U是一个相似映射，证明完成。

以下是不使用超限递归定理的另一种证明思路。 设X$_0$为X中满足以下条件的元素a的集合：存在Y中的元素b，使得s$($a$)$与s$($b$)$相似。对于X$_0$中的每个a，记U$($a$)$为Y中对应的（唯一确定的）b，并令Y$_0$为U的值域。由此可得，要么X$_0$ = X，要么X$_0$是X的一个初始段且Y$_0$ = Y。

习题。良序集 \(X\) 的每个子集要么与 \(X\) 相似，要么与 \(X\) 的某个初始段相似。若 \(X\) 与 \(Y\) 是良序集且 \(X \cong Y\)（即 \(X\) 与 \(Y\) 相似），则该相似映射将 \(X\) 的每个子集的上确界（若存在）映至该子集像集的上确界。


\chapter{序数}

集合\(x\)的后继\(x^+\)定义为\(x \cup \{x\}\)，而\(\omega\)被构造为包含0且每当包含\(x\)时也包含\(x^+\)的最小集合。如果我们从\(\omega\)开始，构造它的后继，再构造该后继的后继，如此无限进行下去，会发生什么？换言之：是否存在某种超越\(\omega, \omega^+, (\omega^+)^+, \ldots\)等的东西，其意义类似于\(\omega\)超越\(0, 1, 2, \ldots\)等？

问题要求一个集合，记为T，包含w，使得T中每个元素（除w本身外）都可以通过反复取后继从w得到。为了更精确地表述这一要求，我们引入一些专门且临时的术语。设一个函数f，其定义域为某个自然数n的严格前驱集合（即dom f = n），若满足f$($0$)$=w（假设n$\neq$0，从而0<n），且当m+<n时f$($m+$)$=$($f$($m$)$)+，则称f为w-后继函数。通过数学归纳法易证，对每个自然数n，存在唯一的定义域为n的w-后继函数。所谓某物要么等于w，要么可通过反复取后继从w得到，即指它属于某个w-后继函数的值域。令S$($n,x$)$表示语句“n是自然数，且x属于定义域为n的w-后继函数的值域”。我们寻找的集合T满足：x$\in$T当且仅当存在某个n使得S$($n,x$)$为真；这样的集合超越w的程度，正如w超越0的程度。

我们知道，对于每个自然数n，我们可以构造集合$\{$x:S$($n,x$)\}$。换句话说，对于每个自然数n，存在一个集合F$($n$)$，使得x$\in$F$($n$)$当且仅当S$($n,x$)$成立。n与F$($n$)$之间的联系看起来非常像函数。然而，事实表明——

然而，我们目前所见的所有集合构造方法都不足以证明存在一个有序对集合F，使得$($n,x$)\in$F当且仅当x$\in$F$($n$)$。为了实现这一显然理想的状态，我们需要再引入一条集合论原理（也是最后一条）。这条新原理大致表述为：对集合的元素进行任何合理的操作，都会产生一个集合。

替换公理。若 S(a,b) 是一个语句，且对于集合 A 中的每个 a，集合 {b:S(a,b)} 均可构成，则存在一个定义域为 A 的函数 F，使得对于 A 中的每个 a，有 F(a)={b:S(a,b)}。

说$\{$b:S$($a,b$)\}$可以形成，当然意味着存在一个集合F$($a$)$，使得b$\in$F$($a$)$当且仅当S$($a,b$)$成立。外延公理表明，替换公理所描述的函数由给定的语句和给定的集合唯一确定。该公理名称的由来在于，它使我们能够通过用新事物替换旧集合中的每个元素，从旧集合中构造出一个新集合。

替换公理的主要应用在于将计数过程远远扩展到自然数之外。从当前观点来看，自然数的关键性质在于它是一个良序集，且每个元素所确定的前段等于该元素本身。（回忆一下，若m和n是自然数，则m<n意味着m$\in$n；这蕴含了$\{$m$\in\omega$:m<n$\}$=n。）扩展计数过程正是基于这一性质；这一思想体系中的基本定义归功于冯$\cdot$诺伊曼。序数被定义为满足对所有$\xi\in\alpha$均有s$(\xi)$=$\xi$的良序集$\alpha$；此处s$(\xi)$如前所述，指前段$\{\eta\in\alpha$:$\eta$<$\xi\}$。

一个不是自然数的序数例子是集合 $\omega$，它由所有自然数组成。这意味着我们现在可以比之前“数”得更远；之前我们唯一可用的数字是 $\omega$ 的元素，而现在我们有了 $\omega$ 本身。我们还有 $\omega$ 的后继 $\omega$⁺；这个集合以明显的方式排序，而且，这个明显的排序是一个满足序数条件的良序。实际上，如果 $\xi\in\omega$⁺，那么根据后继的定义，要么 $\xi\in\omega$（此时我们已经知道 s$(\xi)$ = $\xi$），要么 $\xi$ = $\omega$（此时根据顺序的定义，s$(\xi)$ = $\omega$），因此同样有 s$(\xi)$ = $\xi$。刚才给出的论证相当一般；它证明了如果 $\alpha$ 是一个序数，那么 $\alpha$⁺ 也是一个序数。由此可知，我们的计数过程现在可以延伸到并包括 $\omega$、$\omega$⁺、$(\omega$⁺$)$⁺，以此类推，直至无穷。

此时我们联系到之前关于$\omega$之后情形的讨论。替换公理容易推出，在$\omega$上存在唯一函数F，满足F$($0$)$=$\omega$，且对每个自然数n有F$($n⁺$)$=$($F$($n$)$)⁺。该函数的值域是我们感兴趣的一个集合；而一个更为重要的集合是$\omega$与该函数值域的并集。其原因只有在我们至少初步了解序数算术之后才会明朗，这个并集通常记作$\omega\cdot$2。若借用序数算术的记法，将F$($n$)$写作$\omega$+n，那么我们可以将$\omega\cdot$2描述为：由所有n（n$\in\omega$）和所有$\omega$+n（n$\in\omega$）构成的集合。

现在容易验证$\omega$²是一个序数。当然，这一验证依赖于$\omega$²中序的定义。在此，该定义及其证明均留作练习；我们正式将注意力转向一些一般性论述，其中包含关于$\omega$²的事实作为简单的特例。

集合X中的一种序（偏序或全序）由其初始段唯一确定。换言之，若R和S是X上的序，且对于每个x$\in$X，x的所有R-前驱的集合与所有S-前驱的集合相同，则R与S相同。无论前驱是取严格意义还是非严格意义，这一断言都是显然的。该断言特别适用于良序集。由此特例可推知：若一个集合可能被良序化以使其成为一个序数，则实现这一目标的方式唯一。集合本身决定了使其成为序数的关系必须为何；若该关系满足要求，则该集合是序数，否则不是。称s$(\xi)$=$\xi$意味着$\xi$的前驱必须恰好是$\xi$的元素。因此所讨论的关系就是属于关系。若对于集合$\alpha$中的元素$\xi$和$\eta$，定义$\eta$<$\xi$为$\eta\in\xi$，则结果要么是$\alpha$的一个良序（使得对$\alpha$中每个$\xi$有s$(\xi)$=$\xi$），要么不是；在前一种情况下$\alpha$是序数，在后一种情况下则不是。

我们以提及最初几个序数的名称来结束对序数的初步讨论。在0,1,2,$\dots $之后是$\omega$，而在$\omega$, $\omega$+1, $\omega$+2,$\dots $之后是$\omega\cdot$2。在$\omega\cdot$2+1（即$\omega\cdot$2的后继）之后是$\omega\cdot$2+2，接着是$\omega\cdot$2+3；然后，在这样开始的序列的所有项之后是$\omega\cdot$3。（此时需要再次应用替换公理。）接下来是$\omega\cdot$3+1, $\omega\cdot$3+2, $\omega\cdot$3+3,$\dots $，之后是$\omega\cdot$4。依此类推，我们依次得到$\omega$, $\omega\cdot$2, $\omega\cdot$3, $\omega\cdot$4,$\dots $。应用替换公理可得到在它们之后的东西，其意义与$\omega$在自然数之后相同；那个东西——

第19节 序 数 77

整个过程是 \( \omega^2 \)。之后，整个过程重新开始：\( \omega^2 + 1, \omega^2 + 2, \ldots, \omega^2 + \omega, \omega^2 + \omega + 1, \omega^2 + \omega + 2, \ldots, \omega^2 + \omega 2, \omega^2 + \omega 2 + 1, \ldots, \omega^2 \)

+w$_3$, $\dots $ ,$\omega$²+$\omega^4$, $\dots $ ,$\omega$²$\cdot$2, $\dots $ ,$\omega$²$\cdot$3, $\dots $ ,$\omega$³, $\dots $ ,$\omega^4$, $\dots $ ,$\omega$^$\omega$, $\dots $ ,$\omega$^$($4$)$,

$\dots $，w$($4$($ , $\dots $。这之后的下一个是 $\varepsilon_0$；接着是 $\varepsilon_0$+1，$\varepsilon_0$+2，$\dots $，$\varepsilon_0$+$\omega$，$\dots $，$\varepsilon_0$+$\omega$²，$\dots $，$\varepsilon_0$+$\omega$^$\omega$，$\dots $，$\varepsilon_0\cdot$2，$\dots $，$\varepsilon_0\cdot\omega$，$\dots $，$\varepsilon_0\cdot\omega$^$\omega$，$\dots $，$\varepsilon_0$²，$\dots \dots \dots $


\chapter{序数集合}

序数，按定义，是一种特殊的良序集；我们接下来考察其特殊性质。

最基本的事实是，序数$\alpha$的每个元素同时也是$\alpha$的子集。（换言之，每个序数都是传递集合。）实际上，若$\xi\in\alpha$，则由于s$(\xi)$=$\xi$，可知$\xi$的每个元素都是$\alpha$中$\xi$的前驱，因而特别地，也是$\alpha$的元素。

若$\xi$是序数$\alpha$的一个元素，则如我们所见，$\xi$是$\alpha$的子集，因此$\xi$（相对于从$\alpha$继承的序关系）是一个良序集。断言：$\xi$实际上是一个序数。事实上，若$\eta\in\xi$，则$\eta$在$\xi$中确定的前段与$\eta$在$\alpha$中确定的前段相同；由于后者等于$\eta$，故前者也等于$\eta$。这一结论的另一种表述方式是：序数的每一个前段都是序数。

接下来要注意的是，如果两个序数相似，则它们相等。为证明这一点，假设$\alpha$和$\beta$是序数，且f是从$\alpha$到$\beta$的一个相似映射；我们将证明对于$\alpha$中的每个$\xi$，有f$(\xi)$=$\xi$。证明采用直接的超限归纳法。记S为满足条件的集合。对于$\alpha$中的每个$\xi$，$\alpha$中不属于s$(\xi)$的最小元素就是$\xi$本身。由于f是相似映射，因此$\beta$中不属于s$(\xi)$在f下的像的最小元素是f$(\xi)$。这些论断表明，如果s$(\xi)\subseteq$ S，那么f$(\xi)$和$\xi$是具有相同初始段的序数，从而f$(\xi)$=$\xi$。由此我们证明了当s$(\xi)\subseteq$ S时，必有$\xi\in$ S。超限归纳原理表明S=$\alpha$，从而推出$\alpha$=$\beta$。

如果$\alpha$和$\beta$是序数，那么特别地，它们都是良序集，因此，要么它们相似，要么其中一个与另一个的某个真截段相似。

若$\beta$与$\alpha$的一个初始段相似，则$\beta$与$\alpha$的一个元素相似。由于$\alpha$的每个元素都是序数，由此可知$\beta$是$\alpha$的一个元素，换言之，$\alpha$是$\beta$的延续。我们已知，若$\alpha$与$\beta$是不同的序数，则以下命题成立：

$\beta$C $\alpha$

a是$\beta$的延续，

彼此等价；若它们成立，则可记作

$\beta$<$\alpha$。

我们刚刚证明的是，任意两个序数都是可比较的；也就是说，若 $\alpha$ 和 $\beta$ 是序数，则要么 $\beta$=$\alpha$，要么 $\beta$<$\alpha$，要么 $\alpha$<$\beta$。

上一段的结果可以表述为：每个序数集合都是全序的。实际上，更强的结论成立：每个序数集合都是良序的。假设E是一个非空的序数集合，令a$\alpha$是E中的一个元素。如果对于E中的所有$\beta$都有$\alpha\leq\beta$，那么$\alpha$就是E的第一个元素，结论成立。若非如此，则E中存在某个元素$\beta$使得$\beta$<$\alpha$，即$\beta\in\alpha$；换句话说，$\alpha\cap$E非空。由于$\alpha$是一个良序集，$\alpha\cap$E有第一个元素，记作$\alpha_0$。若$\beta\in$E，则要么$\alpha\leq\beta$（此时$\alpha_0$<$\beta$），要么$\beta$<$\alpha$（此时$\beta\in\alpha\cap$E，因此$\alpha_0\leq\beta$），这就证明了E有第一个元素，即$\alpha_0$。

一些序数是有限的，它们就是自然数（即$\omega$的元素）。其余的被称为超限序数；所有自然数的集合$\omega$是最小的超限序数。每个有限序数（除0外）都有一个直接前驱。如果一个超限序数$\alpha$有一个直接前驱$\beta$，那么，就像自然数一样，$\alpha$=$\beta$⁺。并非每个超限序数都有直接前驱；那些没有直接前驱的被称为极限数。

假设现在e是一个序数的集合。根据我们刚才所见，由于e是一个连续链，因此e中集合的并集a是一个良序集，且对于e中每一个不同于a本身的$\xi$，$\alpha$都是$\xi$的一个连续。由$\alpha$中某个元素所确定的初始段，与该元素在e中任何集合里所确定的初始段相同；这意味着a是一个序数。若$\xi\in$e，则$\xi\leq\alpha$；数a是e中元素的一个上界。

e. 若 $\beta$ 是 e 的另一个上界，则当 $\xi\in$ e 时总有 $\xi$ C $\beta$，因此根据并集的定义，$\alpha$ C $\beta$。这表明 $\alpha$ 是 e 的最小上界；由此我们证明了每一序数集合都有上确界。

是否存在一个恰好由所有序数构成的集合？显然答案是否定的。若存在这样的集合，我们便可构造所有序数的上确界。该上确界将是一个大于或等于每个序数的序数。然而，由于每个序数都存在严格大于它的序数（例如其后继序数），这是不可能的——谈论“所有序数的集合”毫无意义。基于存在这种集合的假设所导致的矛盾，被称为布拉利-福尔蒂悖论（Burali-Forti是一个人，不是两个人）。

我们的下一个目标是证明序数的概念并不像它看起来那样特殊，事实上，每个良序集在所有本质方面都与某个序数相似。这里的“相似”是指技术意义上的同构。该结果的一个非正式表述是：每个良序集都可以被计数。

计数定理。每个良序集都相似于唯一的序数。

证明。由于对于序数而言，相似性等同于相等性，唯一性是显然的。现在假设 \(X\) 是一个良序集，并假设 \(X\) 中的一个元素 \(a\) 满足：由 \(a\) 的每个前驱所确定的初始段都与某个（必然唯一的）序数相似。设 \(S(x, \alpha)\) 表示语句“\(\alpha\) 是一个序数且 \(s(x) \cong \alpha\)”，那么对于 \(s(a)\) 中的每个 \(x\)，可以构造集合 \(\{\alpha : S(x, \alpha)\}\)；事实上，该集合是单元素集。代入公理保证了存在一个集合，它恰好由与 \(a\) 的前驱所确定的初始段相似的序数组成。由此可得，无论 \(a\) 是某个前驱的直接后继，还是所有前驱的上确界，\(s(a)\) 都与某个序数相似。这一论证为应用超限归纳原理铺平了道路；结论是 \(X\) 中的每个初始段都与某个序数相似。

这一事实反过来又证明了替换公理的另一次应用，正如上文所做的那样；最终结论是，如所期望的，X 与某个序数相似。


\chapter{序数算术}

对于自然数，我们曾利用递归定理定义了算术运算，随后证明了这些运算与集合论中的运算以多种理想方式相关联。例如，我们知道两个不相交有限集合E和F的并集元素个数等于\#$($E$)$+\#$($F$)$。现在我们注意到，这一事实本可用于定义加法。若m和n为自然数，我们可以通过寻找不相交集合E和F（满足\#$($E$)$=m且\#$($F$)$=n），并令m+n=\#$($E$\cup$F$)$来定义它们的和。

对应于自然数中已完成和可能完成的操作，序数算术有两种标准方法。部分出于多样性的考虑，部分因为在此语境下递归似乎不太自然，我们将强调集合论方法而非递归方法。

首先指出，将两个良序集合并成一个新的良序集有一种或多或少显而易见的方法。非正式地说，其思路是先写出其中一个集合，然后紧接着写出另一个集合。若要严格表述这一过程，我们立即会遇到一个困难：这两个集合可能并非互不相交。当两个集合存在公共元素时，我们应如何书写这些元素？解决这一困难的方法是使这两个集合互不相交。这可以通过给它们的元素涂上不同颜色来实现。用更数学化的语言来说，就是将每个集合的元素替换为这些元素与某个区分性对象组成的对，并对两个集合使用不同的区分对象。用完全数学化的语言表述：若 \(E\) 和 \(F\) 是任意集合，令 \(E\) 为所有形如 \((x,0)\) 的有序对的集合，其中 \(x\) 属于 \(E\)；

设 \( E \) 为所有形如 \((x,0)\) 的有序对构成的集合，其中 \( x \in F \)。设 \( F \) 为所有形如 \((x,1)\) 的有序对构成的集合，其中 \( x \in F \)。集合 \( E \) 与 \( F \) 显然不相交。存在一个明显的双射对应关系：\( E \) 与 \( E \) 之间（\( x \mapsto (x,0) \)），以及另一个 \( F \) 与 \( F \) 之间（\( x \mapsto (x,1) \)）。

这些对应关系可用于将E和F可能具有的任何结构（例如顺序）传递到 和F。由此可知，无论何时给定两个集合（无论是否带有附加结构），我们总可以用具有相同结构的不相交集合替换它们，因此可以不失一般性地假设它们原本就是不相交的。

在将此构造应用于序数算术之前，我们注意到它可以推广到任意集合族。事实上，若$\{$E$\}$是一个集合族，记Ei为所有有序对$($x,i$)$的集合，其中x属于Ei。（换言之，Ei = Ei $\times\{$i$\}$。）此时族$\{$Ei$\}$是两两不交的，并且它能完成原族$\{$Ei$\}$所能做到的一切。

设 \(E\) 与 \(F\) 为不相交的良序集。在 \(E \cup F\) 上定义序，使得 \(E\) 中元素对以及 \(F\) 中元素对保持原有的顺序，且 \(E\) 中每个元素均排在 \(F\) 中每个元素之前。（用超形式化的语言表述：若 \(R\) 和 \(S\) 分别为 \(E\) 和 \(F\) 上的给定序关系，则令 \(E \cup F\) 按 \(R \cup S \cup (E \times F)\) 排序。）由于 \(E\) 和 \(F\) 是良序的，可知 \(E \cup F\) 也是良序的。良序集 \(E \cup F\) 称为良序集 \(E\) 与 \(F\) 的序数和。

有一种简单且值得推广的方法，可以将序数和的概念扩展到无穷多个被加项。假设 $\{$E$\}$ 是一个由良序集构成的不交族，其索引集 I 是良序的。该族的序数和定义为并集 UiEi，并按如下方式排序：若 a 和 b 是该并集中的元素，且 a$\in$Ei、b$\in$Ej，则 a<b 意味着要么 i<j，要么 i=j 且 a 在 Ei 的给定顺序中先于 b。

序数的加法定义如今已是小儿科。对于每个良序集X，令ord X为与X相似的唯一序数。

X.（若X是有限的，则ord X与之前定义的自然数\#$($X$)$相同。）设$\alpha$和$\beta$为序数，令A和B为不相交的良序集，且ord A=$\alpha$，ord B=$\beta$，并令C为A与B的序数和。根据定义，和$\alpha$+$\beta$即为C的序数，因此ord A+ord B=ord C。需特别注意的是，和$\alpha$+$\beta$与A和B的具体选取无关；任何其他具有相同序数的不相交集对，都会得到相同的结果。

这些考虑可以毫无困难地推广到无穷情形。若 $\{\alpha$i$\}$ 是由良序集 I 索引的序数良序族，设 $\{$Ai$\}$ 为互不相交的良序集族，且对每个 i 有 ord Ai = $\alpha$i，并设 A 为族 $\{$Ai$\}$ 的序数和。则和 $\sum$i ord Ai 按定义即为 A 的序数，故 ord A = $\sum$i $\alpha$i。

=ord A。这里最终结果同样独立于良序集A的任意选择；任何其他选择（具有相同序数）都会给出相同的和。

序数加法的一些性质是好的，另一些则是坏的。在好的方面，有这些恒等式：

以及结合律

$\alpha$+$(\beta$+$\gamma)$=$(\alpha$+$\beta)$+$\gamma$。

同样值得称赞的是，a<$\beta$当且仅当存在一个不同于0的序数$\gamma$使得$\beta$=$\alpha$+$\gamma$。所有这些论断的证明都是初等的。

几乎所有加法的不良行为都源于交换律的失效。例如：1 + $\omega$ = $\omega$（但正如我们刚才所见，$\omega$ + 1 $\neq\omega$）。加法的这种异常行为直观地反映了序关系中的某些事实。例如，若我们在一个无穷序列（类型为 $\omega$）的前端添加一个新元素，结果显然与原先的序列相似；但若将其添加在末尾，则破坏了这种相似性——原集合没有最后一个元素，而新集合却有了一个。

无穷和的主要用途在于激发并促进对乘积的研究。若 \(A\) 与 \(B\) 为良序集，则自然将其乘积定义为将 \(A\) 自身重复相加 \(B\) 次的结果。为赋予此定义以意义，首先需构造一个由良序集构成的不交族，其中每个良序集均与 \(A\) 相似，且以集合 \(B\) 为指标集。实现这一构造的通用方法在此处适用：只需对每个 \(b \in B\) 记 \(A\_b = A \times \{b\}\)。现在，若将序数和的定义应用于族 \(\{A\_b\}\)，则可得出如下定义：两个良序集 \(A\) 与 \(B\) 的序数乘积为笛卡尔积 \(A \times B\) 赋予逆字典序后的结果。换言之，若 \((a,b)\) 与 \((c,d)\) 属于 \(A \times B\)，则 \((a,b) < (c,d)\) 当且仅当 \(b < d\)，或 \(b = d\) 且 \(a < c\)。

若 $\alpha$ 与 $\beta$ 为序数，设 A 与 B 为良序集，且 ord A = $\alpha$，ord B = $\beta$，并设 C 为 A 与 B 的序数积。乘积 $\alpha\beta$ 按定义即为 C 的序数，故有 (ord A$)$(ord B$)$ = ord C。该乘积的定义是明确的，与良序集 A 与 B 的任意选取无关。换言之，在此处我们本可完全避免任何任意性，通过

回忆一下，最容易得到的序数为$\alpha$的良序集就是序数$\alpha$本身（对$\beta$同理）。

与加法类似，乘法也有其优点和缺点。其优点之一是存在单位元。

a0=0, 0$\alpha$=0, $\alpha$1=$\alpha$, l$\alpha$=$\alpha$,

结合律

$\alpha(\beta$Y$)$=$(\alpha\beta)\gamma$

左分配律

$\alpha(\beta$+$\gamma)$=$\alpha\beta$+$\alpha\gamma$

并且，如果两个序数的乘积为零，则其中一个因子必为零。（注意，我们采用乘法优先于加法的标准约定；$\alpha\beta$+$\alpha\gamma$ 表示 ($\alpha\beta)$+$(\alpha\gamma)$。）

乘法交换律不成立，其许多推论也随之失效。例如，2$\omega$ = $\omega$（可视为无穷多个有序对构成的序列），但$\omega$2 $\neq\omega$（可视为一个由无穷序列构成的有序对）。右分配律同样不成立，即$(\alpha$+$\beta)\gamma$通常不等于$\alpha\gamma$+$\beta\gamma$。例如：$($1+1$)\omega$ = 2$\omega$ = $\omega$，但1$\omega$+1$\omega$ = $\omega$+$\omega$ = $\omega$2。

正如重复加法引出了序数乘积的定义，重复乘法可用于定义序数指数。另一种方式是通过超限递归来定义指数运算。其精确细节属于一个广泛且高度专业化的序数理论体系。在此，我们仅简要提示其定义并提及最直接的结果。为定义$\alpha$^$\beta$（其中$\alpha$和$\beta$为序数），采用超限归纳法（对$\beta$进行归纳）。首先规定$\alpha^0$=1，$\alpha$^$(\beta$+1$)$=$\alpha$^$\beta\cdot\alpha$；若$\beta$为极限序数，则定义$\alpha$^$\beta$为所有形如$\alpha$^$\gamma$（$\gamma$<$\beta$）的序数的上确界。若此定义框架被严谨表述，则可推出以下结论。

0 = 0$(\alpha\geq$ 1$)$,

1Y=1，

a³+r=a³a,

a$^8$Y=$($a$^6)$Y

并非所有常见的指数法则都成立；例如，$($a^$\beta)$^$\gamma$ 通常不同于 a^$(\beta\gamma)$。举例：$($2$\cdot$2$)$ = 4^$\omega$ = $\omega$，但 2^0$\cdot$2^4 = $\omega\cdot\omega$ = $\omega$²。

警告：此处及下文使用的序数指数符号与我们之前的用法不一致。从$\omega$到2的所有函数构成的无序集合2，与作为序数序列2, 2$\cdot$2, 2$\cdot$2$\cdot$2等的最小上界的有序集合2，完全是两回事。对此毫无办法；数学用法已在两个阵营中牢固确立。如果在特定情境下，上下文未能揭示应使用两种解释中的哪一种，则必须给出明确的文字说明。


\chapter{施罗德-伯恩斯坦定理}

计数的目的是比较一个集合与另一个集合的大小；最熟悉的计数集合元素的方法是将其按某种适当顺序排列。序数理论是对这一方法的巧妙抽象，但在实现该目的方面略有不足。这并非说序数毫无用处，只是其主要用途在其他领域，例如拓扑学中作为富有启发性的例子和反例的来源。在接下来的内容中，我们仍会关注序数，但它们将不再占据中心地位。（了解我们实际上可以完全舍弃它们这一点颇为重要。基数理论可以在借助序数或不借助序数的情况下构建；两种构建方式各有优势。）在结束这些开场白后，我们转向比较集合大小的问题。

问题在于比较两个集合的大小，当它们的元素看似毫无关联时。判断法国人口多于巴黎人口是容易的，但要比较宇宙的年龄（以秒计）与巴黎的人口（以电子数计）就没那么简单了。对于数学上的例子，考虑以下以辅助集合A定义的两对集合：(i) X = A, Y = A⁺；(ii) X = 𝒫(A), Y = 2^A；(iii) X 是A到自身的所有一一映射的集合，Y 是A的所有有限子集的集合。在每种情况下，我们都可以问集合X和Y哪个有更多元素。问题首先在于为这个问题找到严格的解释，然后解决它。良序定理告诉我们，每个集合都可以被良序化。对于良序集，我们似乎有了一种合理的度量方法。

大小，即它们的序数。这两个评述能解决问题吗？为了比较X和Y的大小，我们能否分别对它们进行良序化，然后比较ord X和ord Y？答案是否定的。问题在于，同一个集合可以有多种良序化方式。一个良序集的序数更多地衡量了该良序本身，而非集合本身。举一个具体例子：考虑所有自然数的集合$\omega$。引入一种新顺序，将0放在所有其他元素之后。（换句话说，如果n和m是非零自然数，则按通常顺序排列；但如果n=0且m$\neq$0，则让m排在n之前。）结果是$\omega$的一种良序化；该良序的序数是$\omega$+1。

若 \(X\) 与 \(Y\) 为良序集，则 \(\text{ord } X < \text{ord } Y\) 的充要条件是 \(X\) 与 \(Y\) 的某个初始段相似。由此可知，即便对序数一无所知，我们仍可比较两个良序集的序大小，只需理解相似性这一概念即可。相似性是为有序集定义的；而对于任意无序集，核心概念则是等价性。（回顾：若存在两个集合 \(X\) 与 \(Y\) 之间的一一对应，则称 \(X\) 与 \(Y\) 等价，记作 \(X \sim Y\)。）若将相似性替换为等价性，则前段所述思路便具有了可操作性。关键在于：若仅需比较大小，我们无需知晓“大小”本身为何物。

若 \(X\) 和 \(Y\) 是集合，且 \(X\) 等价于 \(Y\) 的某个子集，则记作

X $\leq$ Y.

该符号是临时性的，不值得赋予永久名称。然而，在其使用期间，为了方便指代，一个合理的说法是：Y 支配 X。对于某个集合 E 的子集 X 和 Y，所有满足 X$\leq$Y 的有序对 (X,Y$)$ 构成 E 的幂集上的一个关系。该符号恰当地暗示了其所指概念的部分性质。由于该符号让人联想到偏序关系，而偏序关系具有自反性、反对称性和传递性，因此我们可以预期支配关系也具有类似的性质。

自反性和传递性不会造成问题。由于每个集合 \(X\) 都等价于自身的一个子集（即 \(X\) 本身），因此对所有 \(X\) 都有 \(X \leq X\)。如果 \(f\) 是 \(X\) 与 \(Y\) 的一个子集之间的一一对应，且 \(g\) 是 \(Y\) 与 \(Z\) 的一个子集之间的一一对应，那么我们可以将 \(g\) 限制在 \(f\) 的值域上，并将结果与 \(f\) 复合；这样

结论是X等价于Z的一个子集。换言之，如果

若 X $\leq$ Y 且 Y $\leq$ Z，则

有趣的问题是关于反对称性的。如果 X $\leq$ Y 且 Y $\leq$ X，我们能否得出 X = Y？这显然是荒谬的；只要 X 和 Y 是等价的，这些假设就成立，而等价的集合不一定相同。那么，如果我们只知道两个集合中的每一个都与另一个的子集等价，我们能对它们说些什么呢？答案包含在以下著名且重要的结果中。

施罗德-伯恩斯坦定理：若X$\leq$Y且Y$\leq$X，则X～Y。

注记。注意到其逆命题（这恰好是对自反性断言的一个显著加强）从支配关系的定义即可平凡地推出。

证明。设f为从X到Y的一一映射，g为从Y到X的一一映射；问题在于构造X与Y之间的一一对应。为方便起见，可假设集合X与Y无公共元素；若此假设不成立，我们可轻易使其成立，故该假设不失一般性。

我们称X中的元素x是Y中元素f(x)的父代，类似地，Y中的元素y是X中元素g(y)的父代。X中的每个元素x都有一个无穷后代序列，即f(x), g(f(x)), f(g(f(x)))等；同样，Y中的元素y的后代为g(y), f(g(y)), g(f(g(y)))等。这一定义意味着序列中的每一项都是其前面所有项的后代；我们也将称序列中的每一项是其后面所有项的祖先。

对于每个元素（无论是属于X还是Y），必然发生以下三种情况之一。若我们尽可能回溯该元素的祖先，那么要么最终到达X中一个没有父元素的元素（这些“孤儿”恰好是X-g(Y)中的元素），要么最终到达Y中一个没有父元素的元素（即Y-f(X)中的元素），要么该谱系无限回溯。 设Xx为X中起源于X的元素集合（即Xx由X-g(Y)的元素及其在X中的所有后代组成），Xy为X中起源于Y的元素集合（即Xy由Y-f(X)的元素在X中的所有后代组成），X。为X中没有无父祖先的元素集合。 类似地将Y划分为三个集合Yx、Yy和Y。

若 \( x \in X\_x \)，则 \( f(x) \in Y\_x \)，且实际上 \( f \) 在 \( X\_x \) 上的限制是 \( X\_x \) 与 \( Y\_x \) 之间的一一对应。若 \( x \in X\_y \)，则 \( x \) 属于反函数 \( g^{-1} \) 的定义域，且 \( g^{-1}(x) \in Y\_y \)；事实上，

g⁻¹ 到 X$\gamma$ 的限制是 Xy 与 Yy 之间的一一对应。最后，若 x $\in$ X，则 f$($x$)\in$ Y，且 f 到 X 的限制是 X 与 Y 之间的一一对应；或者，若 x $\in$ X，则 g⁻¹$($x$)\in$ Y，且 g⁻¹ 到 X‰ 的限制是 X。与 Y 之间的一一对应。通过组合这三个一一对应，我们得到 X 与 Y 之间的一一对应。

习题。假设 \( f \) 是从 \( X \) 到 \( Y \) 的映射，\( g \) 是从 \( Y \) 到 \( X \) 的映射。证明存在 \( X \) 的子集 \( A \) 和 \( Y \) 的子集 \( B \)，使得 \( f(A) = B \) 且 \( g(Y - B) = X - A \)。这一结果可用于给出 Schröder-Bernstein 定理的一个证明，该证明与上述证明看起来截然不同。

至此我们知道，支配关系具有偏序的基本性质；我们通过指出该序实际上是全序来结束这一初步讨论。这一论断被称为集合的可比性定理：它表明，若X和Y是集合，则要么X$\leq$Y，要么Y$\leq$X。其证明是良序定理与良序集可比性定理的直接推论。将X和Y分别良序化，并利用如下事实：要么这样得到的良序集相似，要么其中一个与另一个的初始段相似；在前一种情形下，X与Y等价，在后一种情形下，其中一个与另一个的子集等价。


\chapter{可数集}

若X与Y为集合，且Y支配X且X支配Y，则施罗德-伯恩斯坦定理适用，表明X与Y等价。若Y支配X但X不支配Y，从而X与Y不等价，此时我们记作

X<Y,

并且我们说Y严格支配X。

支配与严格支配可简洁地表述关于有限集与无限集的一些事实。回顾：若集合X与某个自然数等价，则称X为有限集；否则为无限集。已知若X$\leq$Y且Y有限，则X有限；且$\omega$是无限的（§13）；还知若X无限，则$\omega\leq$X（§15）。最后一条断言的逆命题也成立，可直接证明（利用有限集不能与其真子集等价的特性），或作为施罗德-伯恩斯坦定理的应用。（若$\omega\leq$X，则不可能存在自然数n使得X～n，否则将有$\omega\leq$n，这与$\omega$是无限集矛盾。）

我们刚刚看到，一个集合X是无穷的当且仅当$\omega\leq$ X；接下来我们将证明X是有限的当且仅当X < $\omega$。该证明依赖于严格支配关系的传递性：若X $\leq$ Y且Y $\leq$ Z，且至少其中一个支配关系是严格的，则X < Z。事实上，显然有X $\leq$ Z。若我们假设Z $\leq$ X，则可得Y $\leq$ X且Z $\leq$ Y，从而（由施罗德-伯恩斯坦定理）有X ~ Y且Y ~ Z，这与严格支配的假设矛盾。现在若X是有限的，则X ~ n对某个自然数n成立，且由于$\omega$是无穷的，有n < $\omega$，因此X < $\omega$。

反之，若X < w，则X必为有限集，否则将有w $\leq$ X，从而w < w，这显然矛盾。

一个集合X被称为可数集（或可列集），如果X $\leq\omega$；称为可数无穷集，如果X $\sim\omega$。显然，可数集要么是有限集，要么是可数无穷集。在接下来的内容中，我们的主要目的是证明：对可数集进行许多集合论构造时，结果仍然是可数集。

我们首先观察到，$\omega$的每个子集都是可数的，进而推导出每个可数集的每个子集也都是可数的。这些事实虽然简单，但很有用。

如果 \(f\) 是从 \(\omega\) 到集合 \(X\) 的满射函数，则 \(X\) 是可数的。证明如下：对于每个 \(x \in X\)，集合 \(f^{-1}(\{x\})\) 非空（这里利用了 \(f\) 的满射性质），因此对每个 \(x \in X\)，我们可以找到一个自然数 \(g(x)\) 使得 \(f(g(x)) = x\)。由于函数 \(g\) 是从 \(X\) 到 \(\omega\) 的一一映射，这证明了 \(X \leq \omega\)。关心此类细节的读者可能注意到，该证明使用了选择公理，并且可能想知道是否存在不依赖该公理的替代证明（确实存在）。同样的说明也适用于本节及其后续章节中的其他少数情况，但我们在此不再赘述。

由前一段可知，集合X是可数的当且仅当存在某个可数集到X的满射函数。一个密切相关的结果是：若Y是任意特定的可数无穷集，则非空集合X为可数的充要条件是存在从Y到X的满射函数。

映射 n$\rightarrow$2n 是 $\omega$ 与所有偶数构成的集合 A 之间的一一对应，因此 A 是可数无穷集。这意味着若 X 是可数集，则存在函数 f 将 A 映射到 X 上。类似地，由于映射 n$\rightarrow$2n+1 是 $\omega$ 与所有奇数构成的集合 B 之间的一一对应，因此若 Y 是可数集，则存在函数 g 将 B 映射到 Y 上。函数 h 在 A 上与 f 一致，在 B 上与 g 一致（即当 x$\in$A 时 h$($x$)$=f$($x$)$，当 x$\in$B 时 h$($x$)$=g$($x$)$），它将 $\omega$ 映射到 X$\cup$Y 上。结论：两个可数集的并集是可数集。由此，通过数学归纳法容易证明，有限个可数集的并集是可数集。同样的结果也可通过模仿两个集合情形下的技巧得到；该方法的基础在于：对每个非零自然数 n，存在 $\omega$ 的一族两两不相交的无限子集 $\{$A\_i$\}$ (i<n$)$，使得它们的并集等于 $\omega$。

同样的方法可以用来证明更多结论。断言：存在

存在一个两两不交的无限子集族$\{$An$\}($新$)$，其并集等于$\omega$。证明此结论的一种方法是，通过沿对角线向下计数，将$\omega$的元素排列成无限阵列，如下所示：

抱歉，我无法处理这个请求。请提供需要翻译的英文文本。

10 15 $\cdot$

16 $\dots $

接着考虑这个阵列的行序列。另一种方法是让 A$_0$ 包含 0 和所有奇数，令 A$_1$ 为将 A$_0$ 中每个非零元素加倍后得到的集合，并归纳地令 Aₙ₊$_1$ 为将 Aₙ 中每个元素加倍后得到的集合（n $\geq$ 1）。无论采用哪种方式（还有其他许多方式），细节都很容易补充。 结论：可数无限个可数集合的并集是可数的。证明：给定可数集合族 $\{$Xₙ$\}$（n $\in\omega$），找到函数族 $\{$fₙ$\}$，使得对每个 n，函数 fₙ 将 Aₙ 映射到 Xₙ 上，并定义一个从 $\omega$ 到 $\cup$ₙ Xₙ 上的函数 f，当 k $\in$ Aₙ 时令 f$($k$)$ = fₙ$($k$)$。这一结果与上一段的结果相结合，意味着可数个可数集合的并集总是可数的。

一个有趣且有用的推论是：两个可数集合的笛卡尔积也是可数的。

X$\times$Y = $\cup$\_$\{$y$\in$Y$\}$ (X$\times\{$y$\})$

由于若 \(X\) 是可数的，则对 \(Y\) 中每个固定的 \(y\)，集合 \(X \times \{y\}\) 显然是可数的（利用一一对应 \(x \rightarrow (x,y)\)），结论可由前一段得出。

习题。证明：可数集的所有有限子集构成的集合是可数集。证明：若全序集 \(X\) 的每一个可数子集都是良序的，则 \(X\) 本身也是良序的。

基于前面的讨论，猜测每个集合都是可数的并非不合理。接下来我们将证明事实并非如此；这一否定结果正是基数理论引人入胜之处。

康托尔定理。

每个集合都严格被其幂集支配，或者说，在

换句话说，

X<0(X)

对于所有X。

证明。存在一个从X到P$($X$)$的自然一一映射，即将X的每个元素x与单元素集$\{$x$\}$相关联。该映射的存在性证明了X $\leq$ P$($X$)$；还需证明X与P$($X$)$不等价。

假设 \( f \) 是从 \( X \) 到 \( \wp(X) \) 的一一映射；我们的目的是证明这一假设会导致矛盾。令 \( A = \{ x \in X : x \notin f(x) \} \)；换言之，\( A \) 由 \( X \) 中那些不属于其对应集合的元素组成。由于 \( A \in \wp(X) \) 且 \( f \) 将 \( X \) 映射到 \( \wp(X) \)，故存在 \( X \) 中的元素 \( a \) 使得 \( f(a) = A \)。元素 \( a \) 要么属于集合 \( A \)，要么不属于。若 \( a \in A \)，则根据 \( A \) 的定义，必有 \( a \notin f(a) \)，而由于 \( f(a) = A \)，这不可能成立。若 \( a \notin A \)，则同样根据 \( A \) 的定义，必有 \( a \in f(a) \)，这同样不可能。矛盾由此产生，康托尔定理的证明至此完成。

由于 \(o(X)\) 总是等价于 \(2^X\)（其中 \(2^X\) 是从 \(X\) 到 \(2\) 的所有函数的集合），康托尔定理表明对所有 \(X\) 均有 \(X < 2^X\)。特别地，若将 \(\omega\) 代入 \(X\) 的位置，则可得出所有自然数集合的集合是不可数的（即非可数、不可枚举），等价地，\(2^\omega\) 是不可数的。此处 \(2^\omega\) 是所有由 0 和 1 组成的无穷序列（即从 \(\omega\) 到 \(2\) 的函数）的集合。注意，若将 \(2^\omega\) 理解为序数幂运算的意义，则 \(2^\omega\) 是可数的（事实上 \(2^\omega = \omega\)）。


\chapter{基数算术}

我们研究集合大小比较的一个成果，是定义了一个新概念——基数，并为每个集合X赋予一个基数，记作card X。这些定义使得每个基数a都存在满足card A = a的集合A。我们还将为基数定义一个序关系，通常用$\leq$表示。这些新概念与已有概念之间的联系很容易描述：card X = card Y当且仅当X～Y，而card X < card Y当且仅当X < Y。（若a和b是基数，a < b当然意味着a $\leq$ b但a $\neq$ b。）

基数的定义可以通过几种不同的方式来探讨，每种方式都有其坚定的支持者。为了尽可能维持和谐，并证明该概念的本质属性与具体方法无关，我们将推迟基本构造的讨论。相反，我们将转而研究基数的算术。在此研究过程中，我们将利用上述基数不等式与集合支配性之间的联系；这种对未来的借用已足以满足当前目的。

设 \(a\) 和 \(b\) 为基数，且 \(A\) 与 \(B\) 为不相交的集合，满足 \(\text{card } A = a\) 和 \(\text{card } B = b\)，则按定义记 \(a + b = \text{card}(A \cup B)\)。若 \(C\) 与 \(D\) 为不相交的集合，满足 \(\text{card } C = a\) 和 \(\text{card } D = b\)，则 \(A \sim C\) 且 \(B \sim D\)；由此可得 \(A \cup B \sim C \cup D\)，因此 \(a + b\) 的定义是明确的，与 \(A\) 和 \(B\) 的任意选取无关。如此定义的基数加法满足交换律（\(a + b = b + a\)）和结合律（\(a + (b + c) = (a + b) + c\)）；这些恒等式可直接由并集运算的相应性质推出。

习题。证明：若 a、b、c 和 d 是基数，且满足 a $\leq$ b 和 c $\leq$ d，则 a + c $\leq$ b + d。

为无穷多个加数定义加法并无困难。若 {aᵢ} 是一族基数，且 {Aᵢ} 是相应指标化的一族两两不交的集合，满足对每个 i 有 card Aᵢ = aᵢ，则我们定义：

如前所述，该定义是明确的。

为定义两个基数a和b的乘积ab，我们选取集合A和B，使得card A = a且card B = b，并记ab = card$($A$\times$B$)$。将A和B替换为等势的集合会得到相同的乘积值。另一种定义方式是将a自身相加b次，这涉及构造无穷和，其中指标集I的基数为b，且对每个i$\in$I有ai = a。读者不难验证，这一替代定义与使用笛卡尔积的定义等价。基数乘法满足交换律（ab = ba）和结合律（a$($bc$)$ = (ab$)$c），且乘法对加法满足分配律（a$($b+c$)$ = ab + ac）；其证明是初等的。

习题。证明：若 a、b、c 和 d 是基数，且满足 a $\leq$ b 和 c $\leq$ d，则 ac $\leq$ bd。

对于无限多个因子的乘法定义并无困难。若 {aᵢ} 是一族基数，且 {Aᵢ} 是相应索引的集合族，满足对每个 i 有 card Aᵢ = aᵢ，则我们定义：

该定义是明确的。

习题. 若 $\{$a\_i$\}($i$\in$I$)$ 和 $\{$b\_i$\}($i$\in$I$)$ 是基数族，且对每个 i$\in$I 有 a\_i<b\_i，则

我们可以像从加法过渡到乘法那样，从乘法过渡到指数。对于基数 a 和 b，定义 a^b 最有益的方式是直接给出，但另一种方法是通过重复乘法。在直接定义中，找到集合 A 和 B，使得 card A = a 且 card B = b，并记 a^b = card A^B。另一种方式是将 a^b 定义为“将 a 自乘 b 次”。更精确地说：构造 $\prod$\_$\{$i$\in$I$\}$ a\_i，其中指标集 I 的基数为 b，且每个 a\_i = a。

I. 熟悉的指数法则成立。即，若 a、b 和 c 为基数，则

基数，那么

习题。证明：若a、b和c是基数，且满足a$\leq$b，则a^c$\leq$b^c。证明：若a和b是有限数，均大于1，且c是无限基数，则a^c=b^c。

前述定义及其推论相当直接且毫不令人意外。若将其仅限于有限集合，则所得即为熟悉的有限算术。该主题的新颖之处在于至少有一项为无穷的求和、求积与幂运算中产生。此处"有限"与"无穷"一词的使用具有非常自然的含义：一个基数若为有限集合的基数，则称为有限基数，否则称为无穷基数。

若a和b是基数，且a为有限、b为无限，则a+b=b。

证明过程如下：设A与B为不相交集合，且A等价于某个自然数k，B为无限集；需证A$\cup$B～B。由于$\omega\leq$B，不妨设$\omega\subseteq$B。定义从A$\cup$B到B的映射f如下：f在A上的限制是A与k之间的一一对应；f在$\omega$上的限制由f$($n$)$=n+k（对所有n成立）给出；f在B-$\omega$上的限制为B-$\omega$上的恒等映射。由于该映射构成A$\cup$B与B之间的一一对应，证明完成。

接下来：如果 \$a\$ 是一个无限基数，那么

a+a=a。

证明过程中，设A为集合且card A = a。由于集合A$\times$2是两个与A等势的不交集合（即A$\times\{$0$\}$和A$\times\{$1$\}$）的并集，因此只需证明A$\times$2与A等势即可。我们将采用的方法虽不能完全证明这一点，但已足够接近。其思路是通过利用A的越来越大的子集来逐步逼近所需的一一对应关系的构造。

确切地说，设S为所有满足以下条件的函数f的集合：f的定义域形如X$\times$2，其中X是A的某个子集，且f是X$\times$2与X之间的一一对应。若X是可数

若A的任意可数无限子集X满足X$\times$2～X，则集合族5非空；至少它包含了A的可数无限子集X与X$\times$2之间的一一对应关系。集合族5按包含关系构成偏序集。由于直接验证可知Zorn引理的条件成立，因此F包含一个极大元f，设其值域为X。

断言：A-X 是 有限的。如果 A-X 是 无限的，那么 它 将 包含 一个 可数 无限 集，记作 Y。通过 将 f 与 Y$\times$2 和 Y 之间的 一一对应 相结合，我们 可以 得到 f 的 一个 真扩张，这与 所假设的 极大性 相矛盾。

由于 card X + card X = card X，且 card A = card X + card(A - X)，而 A - X 是有限的，这就完成了 card A + card A = card A 的证明。

以下是基数加法运算的又一结果：若a与b为基数，且至少其中之一为无穷基数，又设c等于a与b中较大者，则

a+b = c.

假设 b 是无穷的，并令 A 与 B 为不相交的集合，且 card A = a，card B = b。由于 a $\leq$ c 且 b $\leq$ c，可得 a + b $\leq$ c + c；又因 c $\leq$ card (A $\cup$ B$)$，故有 c $\leq$ a + b。结果由基数序的反对称性得出。

基数乘法算术中的主要结果是：如果 a 是一个无限基数，那么

a$\cdot$a = a。

该证明与相应的加法事实的证明类似。设 5 为所有函数 f 的集合，其中 f 的定义域形如 X$\times$X，X 是 A 的某个子集，且 f 是 X$\times$X 与 X 之间的一一对应。若 X 是 A 的可数无穷子集，则 X$\times$X～X。这意味着集合 s 非空；至少它包含了 A 中所有可数无穷子集 X 对应的 X$\times$X 与 X 之间的一一对应。集合 S 通过扩展关系构成偏序集。佐恩引理的条件容易验证，由此可知 F 包含一个极大元 f，设其值域为 X。由于 (card X$)$(card X$)$=card X，只需证明 card X=card A 即可完成证明。

假设 card X < card A。由于 card A 等于 card X 与 card(A-X) 中较大的一个，这意味着 card A = card (A-X)，从而推出 card X < card (A-X)。由此可得 A-X

存在一个子集 Y 与 X 等价。由于互不相交的集合 X$\times$Y、Y$\times$X 和 Y$\times$Y 都是无穷集且与 X$\times$X 等价，从而与 X 等价，进而与 Y 等价，因此它们的并集与 Y 等价。将 f 与该并集到 Y 的一一对应相结合，我们得到 f 的一个真扩张，这与所假设的极大性矛盾。这表明我们当前的假设（card X < card A）不成立，从而完成了证明。

习题。证明：若a和b是基数，且至少其中之一为无穷，则a+b=ab。证明：若a和b是基数，且a为无穷而b为有限，则a$^6$=a。


\chapter{基数}

目前我们对基数已有相当了解，但仍不清楚其本质。粗略而言，可称集合的基数是该集合与所有与其等势的集合所共有的性质。我们或许试图通过将X的基数定义为所有与X等势的集合构成的集合来精确化这一概念，但此尝试注定失败——因为不存在如此庞大的集合。受自然数定义方法的启发，下一步自然想到将集合X的基数定义为某个精心挑选的、与X等势的特定集合。这正是我们接下来要做的。

对于每个集合X，存在太多与X等势的其他集合；我们的首要问题是缩小范围。由于我们知道每个集合都与某个序数等势，因此在序数中寻找典型集合、代表集合并非不自然。

可以肯定，一个集合可能与多个序数等势。然而，一个令人鼓舞的事实是：对于每个集合X，与X等势的序数构成一个集合。为证明这一点，首先注意到，很容易找到一个序数，它严格大于所有与X等势的序数。事实上，假设$\gamma$是一个与X的幂集等势的序数。如果$\alpha$是与X等势的序数，那么集合$\alpha$被集合$\gamma$严格支配（即基数$\alpha$ < 基数$\gamma$）。由此可知，不可能有$\gamma\leq\alpha$，因此必须有$\alpha$ < $\gamma$。由于对于序数而言，$\alpha$ < $\gamma$与$\alpha\in\gamma$是同一回事，我们便找到了一个集合，即$\gamma$，它包含了所有与X等势的序数，这意味着与X等势的序数确实构成一个集合。

在与X等价的序数中，哪一个应当被单独挑出并称为X的基数？这个问题只有一个自然的答案。每一组序数都是良序的；

良序集的最小元素似乎是唯一一个需要特别关注的元素。

现在我们准备给出定义：一个基数是一个序数$\alpha$，使得如果$\beta$是与$\alpha$等价的序数（即card $\alpha$ = card $\beta$），则$\alpha\leq\beta$。具有这一性质的序数也称为初始数。若X是一个集合，则card X（X的基数，也称为X的势）是与X等价的最小序数。

习题。证明每一个无限基数都是极限数。

由于每个集合都与其基数等价，因此若 card X = card Y，则 X～Y。反之，若 X～Y，则 card X ~ card Y。由于 card X 是与 X 等价的最小序数，可得 card X $\leq$ card Y；又因 X 与 Y 地位对称，同样有 card Y $\leq$ card X。换言之，card X = card Y 当且仅当 X～Y——这正是基数算术发展中所需的基数条件之一。

一个有限序数（即自然数）不与任何异于自身的有限序数等价。由此可知，若X是有限的，则与X等价的序数集合是单元素集，因此X的基数与X的序数相同。基数和序数都是自然数的推广；在熟悉的有限情形中，这两种推广都与最初产生它们的特例一致。作为这些论述的一个近乎平凡的推论，我们现在可以计算幂集$\wp($A$)$的基数：若card A = a，则card $\wp($A$)$ = 2^a。（注意，这一结果虽然简单，但在此之前无法陈述；直到现在，我们才知道2是一个基数。）证明可直接由$\wp($A$)$与2^A等价这一事实得出。

如果$\alpha$和$\beta$是序数，我们知道$\alpha$<$\beta$或$\alpha\leq\beta$的含义。由此可知，基数自然具备一种序关系。该序满足我们在讨论基数运算时所借用的条件。事实上：若card X<

基数 Y，则 X 是基数 Y 的子集，由此可得 X$\leq$Y。若 X Y，则如前所述，应有 card X = card Y；由此可知必有 X<Y。最后，若 X<Y，则不可能有 card Y $\leq$ card X（因为相似性蕴含等价性），故 card X < card Y。

card Y，则 X 是 card Y 的子集，因此 X $\leq$ Y。

作为这些考虑的一个应用，我们提及该不等式。

对所有基数 \(a\) 成立。 证明：若 \(A\) 是满足 \(\text{card } A = a\) 的集合，则 \(A < \mathcal{P}(A)\)，因此 \(\text{card } A < \text{card } \mathcal{P}(A)\)，从而 \(a < 2^a\)。

习题。若 card A = a，则 A 到自身所有一一映射的集合的基数是多少？A 的所有可数无限子集的集合的基数是多少？

关于序数排序的事实同时也是关于基数排序的事实。例如，我们知道任意两个基数都是可比较的（总是要么 a<b，要么 a=b，要么 b<a），并且实际上每个基数集合都是良序的。我们还知道每个基数集合都有一个上界（实际上是一个上确界），而且，进一步地，对于每个基数集合，都存在一个严格大于其中任何一个基数的基数。这当然意味着不存在最大的基数，或者等价地说，不存在恰好包含所有基数的集合。基于存在这样一个集合的假设所导致的矛盾，被称为康托尔悖论。

基数作为特殊的序数，简化了理论的某些方面，但同时也可能引发一些必须避免的混淆。一个主要的困难来源是算术运算的符号表示。若a和b是基数，则它们同时也是序数，因此和a+b具有两种可能的含义。两个基数的基数之和通常不等于它们的序数之和。这听起来比实际情况更复杂；在实践中，避免混淆并不困难。通过上下文、为基数使用特殊符号，以及偶尔的明确提示，讨论可以相当顺畅地进行。

习题。证明：若$\alpha$和$\beta$是序数，则card$(\alpha$+$\beta)$=card$\alpha$+card$\beta$且card$(\alpha\beta)$=$($card$\alpha)$(card$\beta)$。左侧运算采用序数解释，右侧采用基数解释。

基数符号中有一个非常常用的特殊符号，即希伯来字母的第一个字母（$\aleph$，aleph）。特别地，最小的超限序数，即$\omega$，也是一个基数，作为基数时，它通常用$\aleph_0$表示。

到目前为止，我们明确命名的每一个序数都是可数的。在集合论的许多应用中，最小的不可数序数（通常记为$\Omega$）扮演着重要角色。$\omega$最重要的性质是它是一个无穷良序集。

有序集，其每个初始段都是有限的；相应地，Q 最重要的性质是它是一个不可数无穷的良序集，其每个初始段都是可数的。

最小的不可数序数$\Omega$显然满足基数的定义条件；在其基数角色中，它通常记为N$_1$。等价地，N$_1$可被刻画为严格大于$\aleph_0$的最小基数，换言之，即在基数排序中$\aleph_0$的直接后继。

No与s之间的算术关系是一个关于基数的著名老问题。我们如何通过算术运算从No得到N？目前已知，最基础的步骤（涉及和与积）只会从No回到No本身。我们所知的从No出发并得到更大结果的最简单操作是形成2^No。因此我们知道N $\leq$ 2^No。这个不等式是严格的吗？是否存在一个严格小于2^No的不可数基数？著名的连续统假设猜测答案为否，即N$_1$ = 2^No。目前唯一确知的是，连续统假设与集合论公理是相容的。

对于每个无穷基数 \(a\)，考虑所有严格小于 \(a\) 的无穷基数构成的集合 \(c(a)\)。若 \(a = \omega\)，则 \(c(a) = \varnothing\)；若 \(a = \aleph\_1\)，则 \(c(a) = \{\aleph\_0\}\)。由于 \(c(a)\) 是良序集，它有一个序数，记为 \(\alpha\)。\(a\) 与 \(\alpha\) 之间的联系通常表示为 \(a = \aleph\_\alpha\)。基数 \(a\) 的等价定义可通过超限归纳法给出：根据该定义，*\(a\)（对于 \(\alpha > 0\)）是严格大于所有 \(\beta < \alpha\) 的 \(\aleph\_\beta\) 的最小基数。广义连续统假设是指对于每个序数 \(\alpha\)，有 \(\aleph\_{\alpha+1} = 2^{\aleph\_\alpha}\) 的猜想。


\end{document}